%Version 3.1 December 2024
% See section 11 of the User Manual for version history
%
%%%%%%%%%%%%%%%%%%%%%%%%%%%%%%%%%%%%%%%%%%%%%%%%%%%%%%%%%%%%%%%%%%%%%%
%%                                                                 %%
%% Please do not use \input{...} to include other tex files.       %%
%% Submit your LaTeX manuscript as one .tex document.              %%
%%                                                                 %%
%% All additional figures and files should be attached             %%
%% separately and not embedded in the \TeX\ document itself.       %%
%%                                                                 %%
%%%%%%%%%%%%%%%%%%%%%%%%%%%%%%%%%%%%%%%%%%%%%%%%%%%%%%%%%%%%%%%%%%%%%

%%\documentclass[referee,sn-basic]{sn-jnl}% referee option is meant for double line spacing

%%=======================================================%%
%% to print line numbers in the margin use lineno option %%
%%=======================================================%%

%%\documentclass[lineno,pdflatex,sn-basic]{sn-jnl}% Basic Springer Nature Reference Style/Chemistry Reference Style

%%=========================================================================================%%
%% the documentclass is set to pdflatex as default. You can delete it if not appropriate.  %%
%%=========================================================================================%%

%%\documentclass[sn-basic]{sn-jnl}% Basic Springer Nature Reference Style/Chemistry Reference Style

%%Note: the following reference styles support Namedate and Numbered referencing. By default the style follows the most common style. To switch between the options you can add or remove “Numbered” in the optional parenthesis. 
%%The option is available for: sn-basic.bst, sn-chicago.bst%  
 
%%\documentclass[pdflatex,sn-nature]{sn-jnl}% Style for submissions to Nature Portfolio journals
%%\documentclass[pdflatex,sn-basic]{sn-jnl}% Basic Springer Nature Reference Style/Chemistry Reference Style
\documentclass[pdflatex,sn-mathphys-num]{sn-jnl}% Math and Physical Sciences Numbered Reference Style
%%\documentclass[pdflatex,sn-mathphys-ay]{sn-jnl}% Math and Physical Sciences Author Year Reference Style
%%\documentclass[pdflatex,sn-aps]{sn-jnl}% American Physical Society (APS) Reference Style
%%\documentclass[pdflatex,sn-vancouver-num]{sn-jnl}% Vancouver Numbered Reference Style
%%\documentclass[pdflatex,sn-vancouver-ay]{sn-jnl}% Vancouver Author Year Reference Style
%%\documentclass[pdflatex,sn-apa]{sn-jnl}% APA Reference Style
%%\documentclass[pdflatex,sn-chicago]{sn-jnl}% Chicago-based Humanities Reference Style

%%%% Standard Packages
%%<additional latex packages if required can be included here>

\usepackage{graphicx}%
\usepackage{multirow}%
\usepackage{amsmath,amssymb,amsfonts}%
\usepackage{amsthm}%
\usepackage{mathrsfs}%
\usepackage[title]{appendix}%
\usepackage{xcolor}%
\usepackage{textcomp}%
\usepackage{manyfoot}%
\usepackage{booktabs}%
\usepackage{algorithm}%
\usepackage{algorithmicx}%
\usepackage{algpseudocode}%
\usepackage{listings}%
\usepackage{color}
\usepackage{subcaption}
\usepackage{caption}
\usepackage{float}
\usepackage{array}
\usepackage{siunitx}
\usepackage{makecell}
\usepackage{rotating}
\usepackage{geometry}
\usepackage{float}
\usepackage{threeparttable}
\usepackage{adjustbox}
\usepackage{pdflscape}

\usepackage{siunitx}


%%%%%=============================================================================%%%%
%%%%  Remarks: This template is provided to aid authors with the preparation
%%%%  of original research articles intended for submission to journals published 
%%%%  by Springer Nature. The guidance has been prepared in partnership with 
%%%%  production teams to conform to Springer Nature technical requirements. 
%%%%  Editorial and presentation requirements differ among journal portfolios and 
%%%%  research disciplines. You may find sections in this template are irrelevant 
%%%%  to your work and are empowered to omit any such section if allowed by the 
%%%%  journal you intend to submit to. The submission guidelines and policies 
%%%%  of the journal take precedence. A detailed User Manual is available in the 
%%%%  template package for technical guidance.
%%%%%=============================================================================%%%%

%% as per the requirement new theorem styles can be included as shown below
\theoremstyle{thmstyleone}%
\newtheorem{theorem}{Theorem}%  meant for continuous numbers
%%\newtheorem{theorem}{Theorem}[section]% meant for sectionwise numbers
%% optional argument [theorem] produces theorem numbering sequence instead of independent numbers for Proposition
\newtheorem{proposition}{Proposition}% 
%%\newtheorem{proposition}{Proposition}% to get separate numbers for theorem and proposition etc.

\theoremstyle{thmstyletwo}%
\newtheorem{example}{Example}%
\newtheorem{remark}{Remark}%

\theoremstyle{thmstylethree}%
\newtheorem{definition}{Definition}%

\theoremstyle{plain} % Standard theorem style

\newtheorem{lemma}{Lemma}

\newtheorem{corollary}{Corollary}

\theoremstyle{definition} % Definition style

\newcommand{\hash}{\textsf{hash}}

\theoremstyle{remark} % Remark style


\theoremstyle{plain} % Standard theorem style
\newtheorem{assumption}[theorem]{Assumption}

\raggedbottom
%%\unnumbered% uncomment this for unnumbered level heads


\begin{document}

\title[A Line Graph-Based Metric Learning Framework for Robust Link Prediction in Complex Networks]{A Line Graph-Based Metric Learning Framework for Robust Link Prediction in Complex Networks}

%%=============================================================%%
%% GivenName	-> \fnm{Joergen W.}
%% Particle	-> \spfx{van der} -> surname prefix
%% FamilyName	-> \sur{Ploeg}
%% Suffix	-> \sfx{IV}
%% \author*[1,2]{\fnm{Joergen W.} \spfx{van der} \sur{Ploeg} 
%%  \sfx{IV}}\email{iauthor@gmail.com}
%%=============================================================%%

\author*[1]{\fnm{Adil Imad Eddine} \sur{Hosni}}\email{hosni.adil.emp@gmail.com,hosni.adilimadeddine@emp.edm.dz}

\author[1]{\fnm{Islam} \sur{Baira}}

\author[1]{\fnm{Badis} \sur{Djamaa}}

\author[1]{\fnm{M’hamed} \sur{Mataoui}}

\author[1]{\fnm{Sidhoum Abdellah} \sur{Hamouda}}

\author[1]{\fnm{Hichem} \sur{Merini}}

\author[1]{\fnm{Mohamed Chakib} \sur{Amrani}}


\affil*[1]{\orgdiv{Department of Computer Science}, \orgname{Ecole Militaire Polytechnique}, \orgaddress{\street{Bordj El-Bahri}, \city{Algiers}, \postcode{16046}, \country{Algeria}}}


%%==================================%%
%% Sample for unstructured abstract %%
%%==================================%%

\abstract{
Link prediction is a central task in network analysis, yet many existing methods suffer from two key limitations. 
First, they model edge relationships only indirectly through node-level comparisons. 
Second, they struggle with the class imbalance commonly observed in real-world graphs. 
Reformulating link prediction as node classification on line graphs offers a principled alternative by treating edges as first-class entities, enabling direct relationship modeling. 
However, this formulation introduces computational challenges due to the quadratic growth of the transformed graph.
The scale of the constructed line graphs, reaching hundreds of thousands of nodes and tens of millions of edges, necessitates high-performance computing resources. 
To address these issues, we propose \textbf{LineML}, a scalable framework that transforms the link prediction challenge into a node classification on line graphs. 
LineML combines three complementary components. 
First, a GraphSAGE-based architecture captures node attributes and structural context through multi-hop neighborhood aggregation. 
Second, an adaptive metric-learning module with degree-biased negative sampling refines embeddings according to example difficulty. 
Finally, a pruning and parallel training system mitigates quadratic complexity via spectral pruning and multi-GPU data-parallel execution, enabling efficient optimization and making line graph learning computationally feasible.
We evaluate LineML on 18 benchmark datasets and compare it against 14 state-of-the-art baseline methods.
The method achieves the highest average rank (1.0) on social and biological networks and a best average rank of 2.4 on citation networks. 
Statistical analysis indicates near-complete dominance on social and biological datasets (Cliff's $\delta \geq 0.97$ against all 14 baselines) and substantial improvements on citation networks ($\delta \geq 0.73$ against 12 of 14 methods), with statistically significant gains across all categories (Wilcoxon signed-rank test, $p < 0.01$). 
In addition, spectral pruning reduces training time by up to 4.6$\times$ on the Physics dataset while preserving predictive performance. 
Our distributed implementation attains near-linear scaling, yielding a 13.98$\times$ speedup over line graph baselines on a dual-GPU system.
Finally, the metric learning component remains effective under high negative sampling ratios, maintaining discriminative representations despite increasing class imbalance.
}

\keywords{
Link Prediction,
Line Graph,
Graph Neural Networks,
Metric Learning,
Triplet Loss,
High Performance Computing
}



\maketitle

\section{Introduction}
\label{sec:introduction}


Link prediction is an essential task within data mining and machine learning that anticipates potential connections between entities in a network \cite{32zhang2017weisfeiler}. 
This capability enables applications across diverse domains, such as social recommendation, enhancing knowledge graphs by identifying missing relationships \cite{nickel2015review},  and restoring metabolic pathways in biochemical networks \cite{oyetunde2017boostgapfill}. 
For example, in social media platforms, link prediction algorithms are used to strengthen user engagement by suggesting new connections that align with user interests. 
In biology, predicting protein interactions aids in understanding disease mechanisms, where discovering unknown relationships can reveal potential therapeutic targets \cite{kovacs2019network}. 
Each of these examples highlights the wide-ranging and practical importance of accurate link prediction solutions.

Existing approaches to link prediction have evolved through several methodological paradigms, each addressing specific aspects of the problem while introducing new limitations.
\emph{Heuristic methods} \cite{barabasi1999emergence,adamic2003friends,23zhou2009predicting} employ interpretable similarity metrics but struggle to capture complex structural dependencies and fail to adapt to heterogeneous network patterns. 
\emph{Embedding-based techniques} \cite{koren2009matrix,28ganaie2022ensemble,29tang2015line} learn low-dimensional representations but often treat edges as implicit byproducts of node comparisons, lacking explicit relational modeling. 
\emph{Non-negative matrix factorization approaches} \cite{yao2024deep} offer interpretable dimensionality reduction but face challenges with network sparsity and require specialized extensions for effective link prediction. 
\emph{Graph neural network methods} \cite{14zhang2018link,33cai2020multi,chamberlain2023elph} leverage powerful feature learning capabilities but typically rely on enclosing subgraph extraction, which introduces computational overhead and may miss global relational patterns. 
All these approaches share two fundamental shortcomings: they model edge relationships only indirectly through node-level comparisons, and they struggle with a pronounced class imbalance (positive edges $<<$ negative edges; e.g., $<1\%$ positives) intrinsic to real-world networks

Reformulating link prediction as node classification on line graphs offers a principled alternative. 
This formulation treats edges as first-class entities, enabling direct relationship modeling through edge-to-edge transformations \cite{cai2021line,zhang2023line,liang2023line}. 
However, this formulation introduces substantial computational challenges due to the quadratic growth of the transformed graph.
An original graph with $n$ nodes and $m$ edges yields a line graph with approximately $m$ nodes and up to $\mathcal{O}(m^2)$ edges in dense cases.
Quadratic scaling in the line graph makes standard processing computationally prohibitive for large networks, necessitating high-performance computing solutions. 



Therefore, we introduce \textbf{LineML}, a framework that reformulates link prediction as node classification on line graphs, representing edges as primary entities and enabling direct edge–edge modeling. The framework integrates three complementary innovations. First, a GraphSAGE-based encoder captures node attributes and topological context through multi-hop neighborhood aggregation. Second, an adaptive metric-learning module with degree-biased negative sampling refines embeddings according to example difficulty, dynamically adjusting similarity thresholds. Finally, a scalable pruning and parallelization system mitigates the quadratic growth of line graphs, supported by a distributed high-performance computing (HPC) implementation for efficient large-scale training. The quadratic scaling of line graphs motivates this HPC integration and validates our parallel training strategy, making networks with millions of edges computationally tractable. We evaluate LineML on 18 benchmark datasets spanning citation, social, biological, and infrastructure networks, comparing it against 14 state-of-the-art baselines. Our experiments assess both functional and non-functional performance and include extensive ablation studies to systematically quantify the contribution of each component.


In summary, the contributions of this paper are as follows: 
\begin{itemize} 
\item We introduce LineML, a link prediction framework that approaches link prediction via node classification on line graphs. on line graphs, enabling direct modeling of edge relationships. The framework integrates a GraphSAGE-based encoder for capturing topological context with adaptive metric learning and degree-biased negative sampling to improve discrimination under class imbalance.

\item We propose a scalable pruning and parallel training strategy that mitigates the quadratic growth of line graphs through spectral pruning and distributed data-parallel training, making large-scale line graph learning computationally feasible in practice.

\item We conduct extensive experiments on 18 benchmark datasets spanning citation, social, biological, and infrastructure networks. LineML achieves the best average rank (1.0 in social/biological and 2.4 in citation networks), with statistically significant improvements and large effect sizes (Cliff’s $\delta \geq 0.97$ against 12 of 14 baselines).

\item We analyze the impact of pruning and parallel training on scalability and efficiency, demonstrating substantial reductions in training time while preserving predictive performance, thereby validating the practicality of the proposed HPC-enabled framework.
\end{itemize}



The paper is organized as follows. 
Section~\ref{sec:related_work} reviews the current landscape of link prediction methods. 
Building upon this foundation, Section~\ref{sec:methodology} details the proposed LineML framework, presenting its architecture and theoretical underpinnings. 
To validate the framework, Section~\ref{sec:experiments} presents a comprehensive experimental evaluation that assesses both its predictive accuracy and its computational efficiency. 
Following this, Section~\ref{sec:SysAblation} conducts systematic ablation studies to analyze the contribution of each architectural component. 
These empirical findings are then integrated and examined in Section~\ref{sec:discussion}, which also discusses the framework's broader limitations. 
Finally, Section~\ref{sec:conclusion} synthesizes the key contributions and proposes directions for future research.

\section{Related Work}
\label{sec:related_work}

Link prediction aims to infer missing or future connections between nodes in a graph. 
It is a central problem in network science with applications in social networks, bioinformatics, and recommender systems.
Research on this problem can be organized into five methodological paradigms.
 Following comprehensive surveys and taxonomies \cite{cai2021line,arrar2024comprehensive,ben2025comprehensive}, we categorize existing methods into five groups.


% ==================== Heuristic-based Methods ====================
\subsection{Heuristic-Based Methods}
Heuristic-based methods calculate the likelihood of a link between two nodes using predefined, interpretable similarity scores derived from network topology. 
These methods are often computationally efficient and can be grouped by the scope of topological information they consider.
First-order heuristics, such as \textit{Common Neighbors} and \textit{Preferential Attachment} \cite{barabasi1999emergence}, rely solely on one-hop neighborhood information. 
Second-order heuristics, including the \textit{Adamic–Adar} index \cite{adamic2003friends} and the \textit{Resource Allocation} index \cite{23zhou2009predicting}, extend the scope to two-hop neighbors. 
High-order heuristics like \textit{Katz} index \cite{24katz1953new}, \textit{Rooted PageRank} \cite{25brin2012reprint}, and \textit{SimRank} \cite{26jeh2002simrank} consider global path information, achieving higher accuracy than first-order methods at the cost of greater computational complexity. 
Traditional heuristics excel in dense networks but struggle with the sparsity and complex structural patterns characteristic of real-world graphs.
Recent research has focused on refining these heuristics to overcome their limitations. 
The \textit{Resource Broadcast (RB)} index \cite{liu2024rb} advances the resource allocation concept by modeling link formation as a multi-path broadcast process. 
It defines four distinct transmission paths between nodes, considering broader topological context beyond immediate common neighbors. This bidirectional approach addresses structural details overlooked by classical indices \cite{kong2025general}. 
To tackle the "zero-similarity" problem in sparse networks, the \textit{Basic-Structural Similarity Link Prediction (BSSLP)} method \cite{liu2025bsslp} introduces a hybrid framework. 
It combines a \textit{basic similarity} component, which guarantees a non-zero baseline score for every node pair, with a \textit{structural similarity} component that captures multi-scale local features. This integration proves particularly effective in cold-start and sparse-network scenarios, where traditional heuristics frequently fail.

% ==================== Embedding-Based Approaches ====================
\subsection{Embedding-Based Approaches}
Embedding-based methods learn low-dimensional continuous vector representations of nodes that preserve key structural properties, providing a flexible alternative to fixed heuristic scores. 
Early foundational work includes matrix factorization \cite{koren2009matrix} and stochastic block models for relational learning. 
Following advances in natural language processing, random-walk-based techniques such as \textit{DeepWalk} and \textit{LINE} \cite{28ganaie2022ensemble,29tang2015line} generate node sequences via random walks and employ skip-gram architectures to capture higher-order proximity. 
\textit{Node2Vec} \cite{40grover2016node2vec} extends this paradigm by introducing biased random walks that interpolate between breadth-first and depth-first exploration, enabling greater control over the trade-off between homophily and structural equivalence. 
More recent methods integrate heterogeneous information and adapt to complex real-world settings. 
For instance, Shu and Dai \cite{Shu2023IndustrialKG} combine pre-trained textual encoders (e.g., BERT) with relational GCN modules and a convolutional fusion stage to jointly model node content and graph structure for link prediction in industrial knowledge graphs.
This approach achieves substantial performance gains over baseline methods \cite{Arrar2024Survey}.

% ==================== Deep Learning and GNN-Based Approaches ====================
\subsection{Deep Learning and Graph Neural Network-Based Approaches}
Deep learning techniques, particularly GNNs, have emerged as the dominant paradigm for link prediction, overcoming limitations of handcrafted heuristics and static embeddings. 
A widely adopted strategy reframes link prediction as binary classification on enclosing subgraphs. 
The \textit{SEAL} framework \cite{14zhang2018link} extracts a local subgraph around each candidate link and applies a GNN to classify the subgraph, naturally capturing multi-hop structural context. Subsequent variants enhance \textit{expressivity} through Weisfeiler–Lehman (WL) labeling schemes \cite{32zhang2017weisfeiler} and improve \textit{scalability} via subgraph sketching. 
For example, Chamberlain et al. propose \textit{ELPH} and \textit{BUDDY} \cite{chamberlain2023elph}, which learn compact, differentiable hash sketches that approximate full subgraph information, enabling fast inference without sacrificing accuracy. Other research directions focus on \textit{interpretability}: \textit{PaGE-Link} \cite{Zhang2023PaGE} generates human-readable path explanations for predictions in heterogeneous graphs, while self-explainable GNNs identify influential neighbor subsets to justify link predictions \cite{Zhu2023SelfExplainable}. Additional lines of inquiry include variational approaches, such as the \textit{Graph Auto-Encoder (GAE)} \cite{31kipf2016variational}, which reconstructs the adjacency matrix through a variational autoencoder, and \textit{Neural Relational Inference (NRI)} \cite{41kipf2018neural}, which infers latent interaction graphs from observed dynamics.

The principles of self-supervised and contrastive learning have also been effectively integrated into link prediction pipelines. 
\textit{Deep Graph Infomax (DGI)} \cite{velickovic2018deep} pioneered an unsupervised approach by maximizing mutual information between local node patches and a global graph summary, learning rich, transferable node embeddings. 
Adversarial training frameworks, such as \textit{Adversarial Network Embedding (ANE)} \cite{dai2017adversarial} and \textit{Generative Adversarial Network Embedding (GANE)} \cite{hong2020gane}, employ a generator-discriminator duel to regularize the embedding process, leading to more robust representations of network connectivity. Recent contrastive methods specifically address challenges like class imbalance. 

Link prediction in temporal and evolving graphs constitutes a distinct and active frontier. 
While Transformer and temporal GNN models offer high accuracy, their computational cost has spurred the development of efficient, structure-aware baselines. Li et al. proposed \textit{EAGLE} \cite{Li2025EAGLE}, an efficient model that separately captures short-term influences via time-aware local aggregation and long-term effects via structure-aware temporal PageRank, reporting substantial speedups. 
Gao et al. introduced \textit{HyperEvent} \cite{Gao2025HyperEvent}, a compact continuous-time model that encodes recent structural context for a lightweight transformer classifier, remaining competitive with more complex architectures. 
Other innovations include topological meta-learning: \textit{TMetaNet} \cite{Li2025TMetaNet} exploits persistent homology signatures across graph snapshots to adapt learning over time and improve robustness. 
Contrastive learning has also been adapted for dynamics; Nouranizadeh et al. \cite{Nouranizadeh2024Contrastive} present a contrastive recurrent GNN that learns representations over temporal paths, boosting discrete-time dynamic link prediction accuracy.
While this represents a vital and rapidly advancing area of research, the detailed modeling of temporal dynamics and evolution in graphs is beyond the scope of the present work, which focuses on static link prediction.

% ==================== Non-Negative Matrix Factorization (NMF) Approaches ====================
\subsection{Non-Negative Matrix Factorization Approaches}
Non-negative Matrix Factorization (NMF) provides a powerful, interpretable framework for dimensionality reduction by decomposing a non-negative data matrix into lower-dimensional, non-negative factors \cite{lee1999learning}. 
Its application to network analysis is natural, as the adjacency matrix $\mathbf{A}$ can be factorized as $\mathbf{A} \approx \mathbf{U}\mathbf{V}^T$, where the rows of $\mathbf{U}$ and $\mathbf{V}$ represent latent source and target features of nodes; their dot product scores potential links \cite{wang2017nonnegative}. 
Standard ``shallow'' NMF is limited to linear representations. 
\textit{Deep NMF} overcomes this via a multi-layer decomposition ($\mathbf{A} \approx \mathbf{U}_1 \mathbf{U}_2 \dots \mathbf{U}_L \mathbf{V}^T$) that learns hierarchical, non-linear features of the network topology \cite{trigeorgis2014deep}. 
Pioneering work by Chen et al. \cite{chen2022deep} demonstrated the effectiveness of this deep architecture for link prediction.
Subsequent research has extended Deep NMF to address specific challenges. 
To combat network sparsity, Yao et al. proposed the \textit{EG-DNMF} framework \cite{yao2024deep}, which incorporates an edge generator to hypothesize probable missing links, enriching the input graph before deep factorization. 
For dynamic networks, Li et al. \cite{li2025unified} introduced a unified temporal framework that incorporates graph regularization to enforce smoothness in the latent features across consecutive network snapshots, allowing the model to learn evolutionarily consistent representations for predicting future links.

% ==================== Line Graph Transformation Approaches ====================
\subsection{Line Graph Transformation Approaches}
An increasingly influential strategy that leverages the \emph{line graph} transformation.
 Cai et al. formally introduced this idea with the \textit{Line Graph-based Link Prediction (LGLP)} framework \cite{cai2021line}. 
 Building on this, contrastive learning has been applied in the line graph domain. 
 \textit{Line Graph Contrastive Learning (LGCL)} \cite{zhang2023line} constructs two views (enclosing subgraph and line graph) for each candidate link and maximizes their mutual information to learn robust edge embeddings. 
 Recent works like \textit{DuoLink} \cite{Tola2025DuoLink} further integrate classical heuristic features with learned line graph embeddings, improving expressivity. 
 The approach has also been successfully extended to \textit{weighted} link prediction. 
 Liang et al. \cite{liang2023line} transform subgraphs into line graphs, compute deep edge embeddings via GNNs, and regress the link strength, achieving state-of-the-art results. 
 Complementary methods like \textit{LWP-WL} combine WL labeling with CNNs for weight prediction \cite{Zulaika2022LWPWL}.









\section{Proposed Solution}
\label{sec:methodology}

\begin{table}[!ht]
\centering
\caption{Summary of Mathematical Notations}
\label{tab:notations}
\footnotesize
\begin{tabular}{@{}p{0.20\textwidth}p{0.75\textwidth}@{}}
\toprule
\textbf{Symbol} & \textbf{Description} \\
\midrule

\multicolumn{2}{@{}l}{\textbf{Original Graph}} \\
\midrule
$G = (V, E)$ & Original undirected graph.\\
$n, m$ & Number of vertices and edges: $n = |V|$, $m = |E|$. \\
$E^+, E^-$ & Positive (existing) and negative (non-existent) edges.\\
$\tilde{E}^-$ & Sampled negative edges: $\tilde{E}^- \subset E^-$ with $|\tilde{E}^-| = |E^+|$ \\
$E_{\text{total}}$ & Balanced edge set: $E_{\text{total}} = E^+ \cup \tilde{E}^-$ \\
$d_u$ & Degree of node $u$: $d_u = |\{v: (u,v) \in E\}|$ \\
$\mathbf{x}_u$ & Feature vector of node $u$: $\mathbf{x}_u \in \mathbb{R}^d$ \\
$\ell(u)$ & Discrete label of node $u$: $\ell(u) \in \mathcal{L}$ \\
$N(u)$ & Neighbor set: $N(u) = \{v: (u,v) \in E\}$ \\
$t(e)$ & Timestamp of edge $e$: $t: E \to \mathbb{R}^+$ \\
\addlinespace

\multicolumn{2}{@{}l}{\textbf{Line Graph}} \\
\midrule
$L(G) = (V_L, E_L)$ & Line graph of $G$. \\
$L(G)' = (V_L, E_L')$ & Pruned line graph.\\
$v_e$ & Node in line graph corresponding to edge $e = (u,v)$ in $G$ \\
$\mathbf{x}_{uv}$ & Feature vector: $\mathbf{x}_{uv} = \mathbf{x}_u \odot \mathbf{x}_v \in \mathbb{R}^{d}$ \\
$\ell(u,v)$ & Composite edge label: $\ell(u,v) = (\min(\ell(u), \ell(v)), \max(\ell(u), \ell(v)))$ \\
$\mathbf{X}_L$ & Feature matrix: $\mathbf{X}_L \in \mathbb{R}^{|V_L| \times d}$, $d_{in} = d$ \\
$\bar{d}_L$ & Average degree in line graph: $\bar{d}_L = \frac{\sum_{i=1}^n d_i^2}{m} - 2$ \\
\addlinespace

\multicolumn{2}{@{}l}{\textbf{Sampling Framework}} \\
\midrule
$\mathcal{P}$ & Negative sampling distribution.\\
$\gamma$ & Concentration parameter for CN sampling \\
$\tau$ & Temperature parameter for hard sampling \\
$\phi_\theta(u,v)$ & Model similarity function \\
\addlinespace

\multicolumn{2}{@{}l}{\textbf{Pruning Framework}} \\
\midrule
$\eta$ & Edge preservation ratio: $\eta = |E_L'|/|E_L|$ \\
$\epsilon$ & Spectral approximation tolerance.\\
$\kappa$ & Spectral distortion factor.\\
$L, L'$ & Laplacian matrices of $L(G)$ and $L(G)'$ \\
$\Omega_{uv}, \tilde{\Omega}_{uv}$ & Exact and approximate effective resistance. \\
$\mathbf{Q}$ & Random projection matrix: $\mathbf{Q} \in \mathbb{R}^{r \times |V_L|}$, $r = O(\log|V_L|/\epsilon^2)$ \\
\addlinespace

\multicolumn{2}{@{}l}{\textbf{GNN Architecture}} \\
\midrule
$K$ & Number of GNN layers \\
$\mathbf{h}_v^{(k)}$ & Hidden representation at layer $k$: $\mathbf{h}_v^{(0)} = \mathbf{x}_v$, $\mathbf{h}_v^{(k)} \in \mathbb{R}^{d_k}$. \\
$\mathbf{W}^{(k)}$ & Learnable weight matrix: $\mathbf{W}^{(k)} \in \mathbb{R}^{d_k \times d_{k-1}}$. \\
$\text{AGGREGATE}^{(k)}$ & Aggregation function (mean, max-pooling, or LSTM). \\
$\mathcal{N}(v)$ & Neighborhood of $v$ in $L(G)$: $\mathcal{N}(v) = \{u: (u,v) \in E_L\}$. \\
$\mathbf{z}_v$ & Final embedding: $\mathbf{z}_v = \mathbf{h}_v^{(K)} \in \mathbb{R}^{d_{out}}$. \\
$\mathbf{Z}$ & Embedding matrix: $\mathbf{Z} \in \mathbb{R}^{|V_L| \times d_{out}}$. \\
\addlinespace

\multicolumn{2}{@{}l}{\textbf{Metric Learning}} \\
\midrule
$\mathcal{T}$ & Set of triplets: $\mathcal{T} = \{(a_i, p_i, n_i)\}$ \\
$d(\mathbf{z}_i, \mathbf{z}_j)$ & Euclidean distance: $d(\mathbf{z}_i, \mathbf{z}_j) = \|\mathbf{z}_i - \mathbf{z}_j\|_2$.\\
$\alpha, \beta$ & Base margin ($\alpha \geq 0$) and adaptivity parameter ($\beta \geq 0$).\\
$\mathbf{Z}_{\text{final}}$ & Metric-aware final embeddings. \\
\addlinespace

\multicolumn{2}{@{}l}{\textbf{Link Classification}} \\
\midrule
$M$ & Number of MLP heads in ensemble. \\
$\mathbf{g}_v^{(m,l)}$ & Hidden representation in head $m$, layer $l$. \\
$\mathbf{W}_l^{(m)}, \mathbf{b}_l^{(m)}$ & Weight matrix and bias for head $m$, layer $l$. \\
$\hat{y}_v$ & Prediction from head $m$ and ensemble average.\\
\bottomrule
\end{tabular}
\end{table}


\begin{figure}[!ht]
    \centering
    \includegraphics[width=0.95\textwidth]{figs/fig_archi_1.pdf}
    \caption{LineML Architecture Overview. The \textbf{LineML} framework consists of six core components: 
(1) link sampling to balance classes, 
(2) line graph generation to reformulate the problem, 
(3) optional pruning for computational efficiency, 
(4) GNN-based embedding learning, 
(5) adaptive metric learning for refined representations, and 
(6) link classification for final predictions.}
    \label{fig:LineMLArchi}
\end{figure}




\subsection{Problem Formulation}
\label{subsec:problem_formulation}

Let $G = (V, E)$ be an undirected graph, where $V = \{v_1, v_2, \dots, v_n\}$ denotes the set of $n$ vertices and $E \subseteq V \times V$ represents the set of observed edges. Each vertex $u \in V$ is associated with a feature vector $\mathbf{x}_u$.

\begin{definition}[Link Prediction]
Given a graph $G = (V, E)$, the link prediction task aims to learn a scoring function
$f: V \times V \rightarrow \mathbb{R}$
that estimates the likelihood of a potential edge $(u,v) \notin E$.
Let $E^+ = E$ denote the set of positive (observed) edges and
$E^- = \{(u,v) \notin E \mid u \neq v\}$ denote the set of negative (non-observed) edges.
\end{definition}

Rather than directly performing edge-level prediction, we reformulate link prediction as a node classification problem on the line graph. This transformation allows edges to be treated as first-order entities, enabling standard node-level learning frameworks to be directly applied.

\begin{definition}[Line Graph Transformation]
Given an undirected graph $G = (V, E)$, its line graph $L(G) = (V_L, E_L)$ is defined as
\[
V_L = \{ v_e \mid e \in E \}, \quad
E_L = \{ (v_e, v_{e'}) \mid e=(u,v),\, e'=(u,w) \in E,\, v \neq w \}.
\]
Each node $v_e \in V_L$ corresponds to an edge $e \in E$, and two nodes in $V_L$ are adjacent if their corresponding edges in $G$ share a common endpoint.
\end{definition}

Since the original graph contains only observed edges, the resulting line graph naturally consists solely of positive samples. To enable supervised learning, we construct a training set by uniformly sampling negative edges $\tilde{E}^- \subset E^-$ such that $|\tilde{E}^-| = |E^+|$. We then build an augmented graph $G_{\text{aug}} = (V, E^+ \cup \tilde{E}^-)$ and derive its corresponding line graph $L(G_{\text{aug}})$.

\begin{definition}[Link Prediction as Node Classification]
Given $G = (V, E)$ with vertex features $\{\mathbf{x}_u\}_{u \in V}$ and sampled negative edges $\tilde{E}^-$, we construct the augmented line graph $L(G_{\text{aug}})$ from $E^+ \cup \tilde{E}^-$. The objective is to learn a node classifier
$g: V_L \rightarrow \{0,1\}$ such that
\[
g(v_e) =
\begin{cases}
1, & \text{if } e \in E^+, \\
0, & \text{if } e \in \tilde{E}^-.
\end{cases}
\]
The learned classifier is expected to generalize to unseen nodes in $V_L$, corresponding to previously unobserved edges in the original graph.
\end{definition}

The central challenge in this setting is to accurately infer missing edges under severe supervision imbalance, as real-world graphs provide only positive links while negative links are unobserved. 
This requires carefully defining (i) how negative edges are sampled and in what proportion, 
(ii) how vertex attributes, edge features, and structural information are jointly incorporated, and 
(iii) how an appropriate decision function is learned over candidate links. 
Moreover, transforming edges into nodes via the line graph introduces additional considerations regarding feature construction and label assignment, since node features must be derived from endpoint attributes and edge-level information. 
An effective formulation must therefore address class imbalance, feature integration across multiple graph modalities, and classifier design, while remaining scalable and generalizable to unseen links.


\subsection{Framework Overview}
\label{sec:framework_overview}


The \textbf{LineML} framework consists of six core components illustrated in figure~\ref{fig:LineMLArchi}: 
(1) link sampling to balance classes, 
(2) line graph generation to reformulate the problem, 
(3) optional pruning for computational efficiency, 
(4) GNN-based embedding learning, 
(5) adaptive metric learning for refined representations, and 
(6) link classification for final predictions.

The pipeline begins with link sampling, which takes $G = (V, E^+)$ as input and produces an edge set $E_{\text{total}} = E^+ \cup \tilde{E}^-$ using negative sampling strategies (Section~\ref{subsec:link_sampling}). 
The edge set feeds into line graph generation (Section~\ref{subsec:line_graph_generation}), transforming $G$ into $L(G_{\text{aug}}) = (V_L, E_L)$. 
Node features are constructed via Hadamard product (with optional edge feature integration), and labels via symmetric pairing.
A pruning module (Section~\ref{subsec:linegraph_pruning}) reduces edge density while preserving structural properties using kNN spectral pruning, effective resistance sampling, or degree-based hierarchical pruning.
The pruned line graph undergoes representation learning through a GraphSAGE-based GNN (Section~\ref{sec:gnn_architecture}), producing structural embeddings $\mathbf{Z}$.
These embeddings are refined through adaptive metric learning (Section~\ref{subsec:metric_learning}), which produces metric-aware embeddings $\mathbf{Z}_{\text{final}}$ using an adaptive triplet loss.
Finally, link classification (Section~\ref{sec:link_classifier}) maps $\mathbf{Z}_{\text{final}}$ to predicted probabilities using an ensemble of multilayer perceptrons.
The entire framework is trained end-to-end, ensuring learned representations are both structurally informative and discriminative.

\subsection{Link Sampling Framework}
\label{subsec:link_sampling}
The severe class imbalance in real-world graphs poses significant computational and learning challenges, making it impractical to treat all non-edges as negative samples. We therefore construct a balanced training set $E_{\text{train}} = E^+ \cup \tilde{E}^-$ by uniformly sampling $\tilde{E}^- \subset E^-$ such that $|\tilde{E}^-| = |E^+|$.


Negative sampling addresses both computational and learning challenges in link prediction. 
Beyond merely balancing classes, it enables tractable training by selecting a manageable subset from the vast space of non-edges ($|E^-| \sim O(|V|^2)$). 
More fundamentally, the sampling strategy serves as a critical pedagogical tool.
It shapes what the model learns by determining which negative examples are presented during training, thereby influencing the learned representations and decision boundaries.

The core design choice is the sampling distribution $\mathcal{P}$ that selects $\tilde{E}^-$. 
This distribution is not merely a technical necessity but a deliberate methodological component that encodes inductive biases about which non-edges are most informative for learning.

\begin{definition}[Negative Sampling Distribution]
A sampling distribution $\mathcal{P}: E^- \to [0,1]$ assigns a probability $p_{uv}$ to each non-edge $(u,v) \in E^-$, satisfying the normalization constraint:
\[
\sum_{(u,v) \in E^-} p_{uv} = 1.
\]
This distribution determines which non-edges are presented to the model as negative examples during training.
\end{definition}

The choice of $\mathcal{P}$ involves important trade-offs. 
A uniform distribution over $E^-$, while unbiased, may include many trivial negatives that offer little learning signal. 
Conversely, overly aggressive sampling of challenging negatives can initially destabilize training if the model lacks capacity to distinguish them. 
Ideally, the sampling strategy should evolve with the model's capability, presenting increasingly difficult examples as learning progresses.
We propose and analyze four complementary sampling strategies, each with distinct mathematical properties and learning implications. 


\subsubsection{Random Negative Sampling}

Random sampling serves as a baseline that draws negative examples uniformly from all non-edges.

\begin{definition}[Uniform Random Sampling]
For any non-edge $(u,v) \in E^-$, the sampling probability is:
\begin{equation}
   p_{uv}^{\text{rand}} = \frac{1}{|E^-|} \approx \frac{1}{|V|^2 - |E|}. 
\end{equation}


\end{definition}

This strategy is unbiased and preserves the global distribution of non-edges. Its key property is that it does not favor any particular structural pattern, providing a "neutral" set of negatives.
Random sampling is computationally efficient with $O(1)$ sampling cost after $O(|E^-|)$ precomputation, but may include many trivial negatives that are easily distinguished from positives.

\subsubsection{Degree-Based Negative Sampling}

Scale-free degree distributions are a hallmark of real-world networks, where a small fraction of high-degree nodes (hubs) account for a disproportionate share of connections. Degree-based negative sampling leverages this heterogeneity to generate informative training examples. Rather than selecting negatives uniformly at random, we bias the sampling toward connections that are statistically plausible yet absent, creating a more challenging and realistic learning signal.

The core mechanism assigns sampling probability proportional to the product of endpoint degrees:
\begin{equation}
    p_{uv}^{\text{deg}} = \frac{d_u \cdot d_v}{\sum_{(x,y) \in E^-} d_x d_y}, \quad \text{for } (u,v) \in E^-,
\end{equation}


where $d_u$ is the degree of node $u$. This formulation embodies a simple but powerful principle: non-edges between high-degree nodes are more informative negatives because their absence is more surprising under a random connection model.

\begin{proposition}[Degree-Consistent Sampling]
Let $d_u^{\text{neg}}$ denote the number of sampled negative edges incident to node $u$. Under degree-product sampling,
\begin{equation}
  \mathbb{E}[d_u^{\text{neg}}] = \frac{d_u}{S} \sum_{v: (u,v) \notin E} d_v, \quad \text{where } S = \sum_{(x,y) \in E^-} d_x d_y.  
\end{equation}


Thus, high-degree nodes participate in more negative examples, preserving the network's inherent heterogeneity in the training data.
\end{proposition}

\begin{proof}
By linearity of expectation,
\begin{equation}
\mathbb{E}[d_u^{\text{neg}}] = \sum_{v: (u,v) \notin E} p_{uv}^{\text{deg}} = \sum_{v: (u,v) \notin E} \frac{d_u d_v}{S} = \frac{d_u}{S} \sum_{v: (u,v) \notin E} d_v.
\end{equation}
The summation $\sum_{v: (u,v) \notin E} d_v$ tends to correlate positively with $d_u$ in networks where high-degree nodes have many potential connections, establishing the approximate proportionality.
\end{proof}

This strategy offers both pedagogical and practical advantages. 
Pedagogically, by presenting negatives that involve hubs the model must learn more than simple degree-based heuristics.
It must distinguish why certain high-degree pairs remain unconnected despite statistical likelihood. 
Practically, degree-based sampling is highly efficient: after computing node degrees ($O(|V|)$), samples can be drawn in $O(\log |V|)$ time using the alias method, requiring only $O(|V|)$ storage. 
This combination of structural realism and computational scalability makes it particularly suitable for large-scale applications where both learning effectiveness and efficiency are paramount.

\subsubsection{Common Neighbor Sampling}

This approach focuses on local topology by favoring non-edges between nodes that share neighbors, capturing the intuition that missing links often occur in locally dense regions.

\begin{definition}[Common Neighbor Sampling]
Let $N(u)$ denote the neighbor set of node $u$. The Jaccard similarity between nodes $u$ and $v$ is:
\begin{equation}
\text{Jaccard}(u,v) = \frac{|N(u) \cap N(v)|}{|N(u) \cup N(v)|}.
\end{equation}
The sampling probability is:
\begin{equation}
p_{uv}^{\text{cn}} = \frac{\text{Jaccard}(u,v)^\gamma}{\sum_{(x,y) \in E^-} \text{Jaccard}(x,y)^\gamma},
\end{equation}
where $\gamma \geq 0$ controls concentration ($\gamma=1$ for linear weighting, $\gamma>1$ for more aggressive sampling of high-similarity pairs).
\end{definition}


This strategy generates challenging negatives that are locally plausible, as they connect nodes that already share social or functional context. It aligns with sociological principles like triadic closure.
but it is not suitible for non social network such infrastructure networks or biology networks

\subsubsection{Hard Negative Sampling}

Hard negative sampling selects examples that are most challenging for the current model, forcing it to learn finer distinctions. This adaptive strategy evolves during training as the model improves.

\begin{definition}[Hard Negative Sampling]
Let $\phi_\theta(u,v)$ measure the similarity between nodes $u$ and $v$ according to the current model parameters $\theta$ (e.g., cosine similarity of their embeddings). Hard negatives are selected with probability:
\begin{equation}
p_{uv}^{\text{hard}} = \frac{\exp(\phi_\theta(u,v)/\tau)}{\sum_{(x,y) \in E^-} \exp(\phi_\theta(x,y)/\tau)},
\end{equation}
where $\tau > 0$ is a temperature parameter. Lower $\tau$ increases the "hardness" by focusing probability mass on the most similar non-edges.
\end{definition}



Hard sampling generates the most informative negatives for improving the current model but requires updating as embeddings change, making it computationally more expensive than static strategies.

\subsection{Line Graph Generation}
\label{subsec:line_graph_generation}

This subsection details the line graph transformation process through five interconnected components: 
(1) formal definition and basic properties, 
(2) transformation of node and edge attributes and labels, 
(3) prevention of data leakage in temporal settings, 
(4) the complete generation algorithm with complexity analysis, and 
(5) theoretical analysis of the resulting structure's properties.

\subsubsection{Definition and Structural Properties}

The transformation of the original graph $G = (V, E_{\text{total}})$ into its line graph representation $L(G)$ represents a fundamental operation of LineML framework.
This transformation converts edges in $G$ to nodes in $L(G)$, thereby allowing the application of powerful node-based graph learning techniques to edge prediction problems. 
The graph $G = (V, E_{\text{total}})$ comprises both positive ($E^+$) and negative ($E^-$) edges obtained through the sampling strategies described in Section~\ref{subsec:link_sampling}.

\begin{theorem}[Line Graph Size Properties]
Given a graph $G = (V, E)$ with $n = |V|$ nodes, $m = |E|$ edges, and degree sequence $\{d_1, d_2, \ldots, d_n\}$, the line graph $L(G)$ has:
\begin{equation}
 |V_L| = m, \quad |E_L| = \frac{1}{2}\left(\sum_{i=1}^{n} d_i^2 - 2m\right) = \sum_{i=1}^{n} \binom{d_i}{2}.
\end{equation}
Moreover, the average degree $\bar{d}_L$ in $L(G)$ satisfies:
\begin{equation}
\bar{d}_L = \frac{\sum_{i=1}^{n} d_i^2}{m} - 2.
\end{equation}
\end{theorem}

\begin{proof}
The equality $|V_L| = m$ follows directly from the bijection between edges in $G$ and nodes in $L(G)$. For the edge count, consider any node $u$ in $G$ with degree $d_u$. Each pair of distinct edges incident to $u$ contributes one edge in $L(G)$. Since there are $\binom{d_u}{2}$ such pairs, summing over all nodes gives $|E_L| = \sum_{i=1}^{n} \binom{d_i}{2}$. Using the identity $\binom{d}{2} = \frac{d^2 - d}{2}$, we have:
\begin{equation}
|E_L| = \frac{1}{2}\left(\sum_{i=1}^{n} d_i^2 - \sum_{i=1}^{n} d_i\right) = \frac{1}{2}\left(\sum_{i=1}^{n} d_i^2 - 2m\right),
\end{equation}
where the last equality follows from the handshaking lemma $\sum_{i=1}^{n} d_i = 2m$. The average degree in $L(G)$ is given by:
\begin{equation}
\bar{d}_L = \frac{2|E_L|}{|V_L|} = \frac{\sum_{i=1}^{n} d_i^2 - 2m}{m} = \frac{\sum_{i=1}^{n} d_i^2}{m} - 2.
\end{equation}
\end{proof}

\begin{corollary}[Density Implications]
For a graph $G$ with average degree $\bar d$, the line graph $L(G)$ satisfies
\begin{equation}
\bar d_L \ge 2\bar d - 2,
\end{equation}
with equality if and only if $G$ is regular. Moreover, for scale-free networks
with power-law degree distribution $P(d)\propto d^{-\gamma}$ $(2<\gamma<3)$,
the quantity $\sum_i d_i^2/m$ diverges with network size, making $L(G)$
increasingly dense.
\end{corollary}



\subsubsection{Node Attribute and Label Transformation}
\label{subsec:attribute_transformation}

Having established the structural properties of the line graph transformation, we now address how to transfer semantic information from the original graph to the line graph representation. This transformation is essential for enabling the GNN to learn meaningful patterns in the edge classification task. 

\paragraph{Symmetric Label Transformation}
    
For each edge $(u, v) \in E$ in the original graph $G$, where nodes $u$ and $v$ possess discrete labels $\ell(u)$ and $\ell(v)$ (representing categorical attributes or community assignments), we construct a composite label for the corresponding line graph node $v_{uv}$. To ensure invariance to edge direction and maintain consistency, we adopt a symmetric labeling scheme inspired by \cite{cai2021line}:
    
\begin{equation}
\ell(u,v) = \left( \min(\ell(u), \ell(v)), \max(\ell(u), \ell(v)) \right),
\label{eq:sym_label}
\end{equation}
    
where the tuple $(a, b)$ with $a \leq b$ uniquely identifies the unordered pair $\{\ell(u), \ell(v)\}$. This transformation has several desirable properties:
    
\begin{theorem}[Properties of Symmetric Edge Labeling]
\label{thm:symmetric_labeling}
Let $G = (V, E)$ be an undirected graph equipped with a vertex labeling function $\ell: V \to \mathcal{C}$, where $\mathcal{C}$ is a discrete label space. Define the induced edge labeling function $\Lambda: E \to \mathcal{C} \times \mathcal{C}$ as:
\[
\Lambda(u, v) = \bigl(\min(\ell(u), \ell(v)), \max(\ell(u), \ell(v))\bigr),
\]
for all $(u, v) \in E$. Then $\Lambda$ satisfies the following properties:
    
\begin{enumerate}
    \item \textbf{Symmetry}: $\Lambda(u, v) = \Lambda(v, u)$ for all $(u, v) \in E$.
    \item \textbf{Label Injectivity}: For any fixed $u \in V$, if $\Lambda(u, v) = \Lambda(u, w)$ for $(u, v), (u, w) \in E$, then $\ell(v) = \ell(w)$.
    \item \textbf{Order Invariance}: $\Lambda(u, v)$ is invariant under permutation of $u$ and $v$.
    \item \textbf{Automorphism Equivariance}: For any graph automorphism $\phi: V \to V$ that preserves vertex labels (i.e., $\ell(\phi(u)) = \ell(u)$ for all $u \in V$), we have $\Lambda(\phi(u), \phi(v)) = \Lambda(u, v)$ for all $(u, v) \in E$.
\end{enumerate}
\end{theorem}
    
\begin{proof}
We prove each property in turn.
    
\textbf{(i) Symmetry:} By definition, $\Lambda(u, v) = (\min(\ell(u), \ell(v)), \max(\ell(u), \ell(v)))$. Since the $\min$ and $\max$ operations are commutative, swapping $u$ and $v$ yields $(\min(\ell(v), \ell(u)), \max(\ell(v), \ell(u))) = \Lambda(v, u)$.

\textbf{(ii) Label Injectivity:} Fix $u \in V$ and let $a = \ell(u)$. Consider $v, w \in V$ such that $\Lambda(u, v) = \Lambda(u, w)$. Denote $\ell(v) = b$ and $\ell(w) = c$. We consider three cases:

\begin{itemize}
    \item If $a \leq b$ and $a \leq c$, then $\Lambda(u, v) = (a, b)$ and $\Lambda(u, w) = (a, c)$. Equality implies $b = c$.
    \item If $b < a$ and $c < a$, then $\Lambda(u, v) = (b, a)$ and $\Lambda(u, w) = (c, a)$. Equality implies $b = c$.
    \item If $a \leq b$ and $c < a$, then $\Lambda(u, v) = (a, b)$ and $\Lambda(u, w) = (c, a)$. Since $a \leq b$ and $c < a$, we have $c < a \leq b$, so the tuples cannot be equal. This case contradicts the assumption.
\end{itemize}

Thus, in all valid cases, $b = c$, hence $\ell(v) = \ell(w)$.

\textbf{(iii) Order Invariance:} This follows directly from the symmetry property, as $\Lambda(u, v) = \Lambda(v, u)$ for any edge $(u, v) \in E$.

\textbf{(iv) Automorphism Equivariance:} Let $\phi$ be an automorphism with $\ell \circ \phi = \ell$. Then for any $(u, v) \in E$:
\begin{align*}
\Lambda(\phi(u), \phi(v)) 
&= \bigl(\min(\ell(\phi(u)), \ell(\phi(v))), \max(\ell(\phi(u)), \ell(\phi(v)))\bigr) \\
&= \bigl(\min(\ell(u), \ell(v)), \max(\ell(u), \ell(v))\bigr) \\
&= \Lambda(u, v).
\end{align*}
Thus, $\Lambda$ is equivariant under label-preserving automorphisms.
\end{proof}

\paragraph{Feature Transformation via Hadamard Product with Edge Feature Integration}

For continuous node attributes, we propose a transformation that combines continuous feature interactions with discrete label information. Let $\mathbf{x}_u \in \mathbb{R}^d$ and $\mathbf{x}_v \in \mathbb{R}^d$ be the attribute vectors of nodes $u$ and $v$ respectively. When edge features $\mathbf{e}_{uv} \in \mathbb{R}^{d_e}$ are available, we integrate them using a gating mechanism inspired by \cite{you2020design}. The feature vector for the corresponding line graph node $v_{uv}$ is constructed as:

\begin{equation}
\mathbf{X}_{uv} = \text{Concat}\left( 
    \left( \mathbf{x}_u \odot \mathbf{x}_v \right) \odot \sigma\left(\mathbf{W}_e \mathbf{e}_{uv} + \mathbf{b}_e\right) + \text{MLP}\left(\mathbf{e}_{uv}\right),     \ell(u,v)\right) \in \mathbb{R}^{d + 2},
\label{eq:attribute_transform_with_edge}
\end{equation}

where $\odot$ denotes the Hadamard (element-wise) product, $\sigma$ is the sigmoid activation function, $\mathbf{W}_e \in \mathbb{R}^{d \times d_e}$ is a learnable projection matrix, $\mathbf{b}_e \in \mathbb{R}^d$ is a bias term, $\text{MLP}$ is a multi-layer perceptron that processes edge features when $d_e > 0$, and $\ell(u,v)$ is the symmetric edge label tuple defined in Equation \ref{eq:sym_label}.

When edge features are unavailable, the transformation simplifies to:

\begin{equation}
\mathbf{X}_{uv} = \text{Concat}\left( 
    \mathbf{x}_u \odot \mathbf{x}_v,
    \ \ell(u,v)
\right) \in \mathbb{R}^{d + 2}.
\label{eq:attribute_transform_no_edge}
\end{equation}

This unified design offers several advantages. The Hadamard product $\mathbf{x}_u \odot \mathbf{x}_v$ captures pairwise multiplicative interactions without increasing dimensionality, providing more expressive edge representations than additive operators while remaining parameter-efficient. Edge features are incorporated via a gating term $\sigma(\mathbf{W}_e \mathbf{e}_{uv} + \mathbf{b}_e)$, enabling adaptive, edge-aware interactions that handle heterogeneous connections. The discrete label tuple $\ell(u,v)$ preserves structural semantics and community information as a compact two-dimensional representation, while the concatenation operation ensures both continuous and discrete information streams are preserved. The formulation is permutation-equivariant for symmetric edge features, ensuring consistent representations for undirected edges. Finally, the operation is computationally efficient, running in $\mathcal{O}(d + d_e)$ time and space, avoiding the dimensional expansion of concatenation while retaining strong representational capacity.

The choice of Hadamard product with edge gating is further justified by its successful application in state-of-the-art graph representation learning. \cite{zhang2021nbfnet} demonstrate that multiplicative interactions between node features achieve superior performance for link prediction tasks compared to additive aggregations. Moreover, \cite{dimchev2023hadamard} provide theoretical analysis showing that Hadamard products in graph networks can approximate any continuous permutation-invariant function on graphs, making them universal approximators while maintaining efficiency.

When edge features are available, our transformation follows the design principles outlined in \cite{you2020design} for incorporating edge information in graph neural networks. The gating mechanism is particularly effective as it allows the model to learn which node feature interactions are relevant given specific edge contexts, implementing a form of feature-wise modulation similar to FiLM layers in computer vision \cite{perez2018film}.
    
\subsubsection{Data Leakage Prevention in Line Graph Construction}
\label{subsec:leakage_prevention}

While attribute transformation ensures feature representation fidelity, another critical consideration is preventing information leakage between temporal splits. 
This challenge is exacerbated by the line graph's adjacency structure, as edges in $L(G)$ connect nodes that correspond to edges in $G$ sharing common endpoints. This creates potential leakage paths if not properly managed.

\begin{definition}[Temporal Partitioning and Causal Information Flow]
\label{def:temporal_partitioning}
Let $t(e)$ denote the timestamp of an edge $e \in E$ (or an arbitrary ordering for
static graphs). The edge set is partitioned into three disjoint temporal folds:
\begin{equation}
E = E_{\text{train}} \cup E_{\text{val}} \cup E_{\text{test}},
\quad
\max t(E_{\text{train}}) < \min t(E_{\text{val}}) < \min t(E_{\text{test}}).
\end{equation}

Let $L = L(G[E])$ be the line graph constructed from the full edge set $E$.
Each node $v_e \in V(L)$ inherits a temporal label
$\tau(v_e) \in \{\text{train}, \text{val}, \text{test}\}$
according to the membership of $e$ in $E_{\text{train}}, E_{\text{val}}$, or
$E_{\text{test}}$.

Information propagation in $L$ is constrained by a causal mask: when computing
the representation of a node $v_e$, only neighboring nodes $v_{e'}$ satisfying
$t(e') \le t(e)$ are allowed to contribute.
\end{definition}

\begin{theorem}[Leakage-Free Information Flow]
\label{thm:leakage_free}
Under the temporal partitioning and causal information-flow constraint in
Definition~\ref{def:temporal_partitioning}, the following properties hold:

\begin{enumerate}
    \item Nodes labeled as $\emph{train}$ receive information only from other
    $\emph{train}$ nodes.

    \item Nodes labeled as $\emph{val}$ may receive information from
    $\emph{train}$ and $\emph{val}$ nodes, but never from $\emph{test}$ nodes.

    \item Nodes labeled as $\emph{test}$ may receive information from
    $\emph{train}$, $\emph{val}$, and $\emph{test}$ nodes.

    \item Consequently, no information from future edges influences
    representations or predictions of past edges.
\end{enumerate}
\end{theorem}


\begin{proof}
In the line graph $L$, two nodes $v_{e_1}$ and $v_{e_2}$ are adjacent if and only
if the corresponding edges $e_1$ and $e_2$ share an endpoint in $G$.
This adjacency is independent of temporal labels.

Leakage prevention is enforced not by removing adjacency, but by restricting
information flow using the causal mask defined in
Definition~\ref{def:temporal_partitioning}.
Specifically, when computing the representation of a node $v_e$, only neighbors
$v_{e'}$ satisfying $t(e') \le t(e)$ are permitted to contribute.

We consider each temporal fold:

\textbf{Train nodes.}
For $e \in E_{\text{train}}$, all edges in $E_{\text{val}} \cup E_{\text{test}}$
satisfy $t(e') > t(e)$ by construction.
Hence, train nodes can receive information only from other train nodes.

\textbf{Validation nodes.}
For $e \in E_{\text{val}}$, edges in $E_{\text{train}}$ and $E_{\text{val}}$
satisfy $t(e') \le t(e)$, while edges in $E_{\text{test}}$ satisfy
$t(e') > t(e)$.
Thus, validation nodes may receive information from train and validation nodes,
but not from test nodes.

\textbf{Test nodes.}
For $e \in E_{\text{test}}$, all edges in
$E_{\text{train}} \cup E_{\text{val}} \cup E_{\text{test}}$
satisfy $t(e') \le t(e)$.
Therefore, test nodes may receive information from all temporal folds.

In all cases, information flows only from past to future, and never in the
reverse direction. This guarantees the absence of temporal leakage.
\end{proof}


\subsubsection{Generation Algorithm and Complexity Analysis}
\label{subsec:algorithm_complexity}

Integrating these principles, we present the complete line graph generation procedure in Algorithm~\ref{algo:line_graph_generation}. The algorithm implements the theoretical principles defined in the preceding sections: it enforces the Temporal Partitioning Constraint, applies the Hadamard product feature transformation  for attributed graphs, and computes the symmetric composite label when discrete node labels are available.

\begin{algorithm}[H]
\caption{Line Graph Generation }
\label{algo:line_graph_generation}
\begin{algorithmic}[1]
\Require Original graph $G = (V, E)$, temporal ordering $t: E \to \mathbb{R}^+$, split ratios $(\eta_{\text{train}}, \eta_{\text{val}}, \eta_{\text{test}})$
\Require \textbf{Optional:} Node features $\{\mathbf{x}_u \in \mathbb{R}^d : u \in V\}$, edge features $\{\mathbf{e}_{uv} \in \mathbb{R}^{d_e} : (u,v) \in E\}$, and/or labels $\{\ell_u : u \in V\}$
\Ensure Line graphs $L_{\text{train}}, L_{\text{val}}, L_{\text{test}}$ with features and labels, with no leakage

\State Sort edges $E$ by $t(e)$ in ascending order
\State Partition: $E_{\text{train}} \gets \text{first } \eta_{\text{train}}|E| \text{ edges}$
\State $E_{\text{val}} \gets \text{next } \eta_{\text{val}}|E| \text{ edges}$
\State $E_{\text{test}} \gets \text{remaining } \eta_{\text{test}}|E| \text{ edges}$

\For{$S \in \{\text{train}, \text{val}, \text{test}\}$}
    \State Initialize $V_L^{(S)} \gets \emptyset$, $E_L^{(S)} \gets \emptyset$
    \State Initialize feature matrix $\mathbf{X}_L^{(S)} \gets [\ ]$, target vector $\mathbf{Y}^{(S)} \gets [\ ]$
    
    \For{each edge $e = (u,v) \in E_S$}
        \State Create node $v_e \in V_L^{(S)}$

            \State $\mathbf{x}_{uv} \gets \mathbf{x}_{uv} \odot \sigma(\mathbf{W}_e \mathbf{e}_{uv} + \mathbf{b}_e) + \text{MLP}(\mathbf{e}_{uv})$ \Comment{Optional edge gating}
      
        \If{node labels available}
            \State $\mathbf{x}_{uv} \gets \text{Concat}(\mathbf{x}_{uv}, (\min(\ell(u),\ell(v)), \max(\ell(u),\ell(v))))$ 
        \EndIf
        
        \State Append $\mathbf{x}_{uv}$ to $\mathbf{X}_L^{(S)}$
  
        \State $y_e \gets \mathbb{I}(e \in E^+)$ \Comment{1 for positive edge, 0 for negative}
        \State Append $y_e$ to $\mathbf{Y}^{(S)}$
    \EndFor
    

    \For{each node $u \in V$}
        \State Let $E_u = \{(u,v) \in E_S\}$ be edges in $E_S$ incident to $u$
        \For{each pair $(u,v), (u,w) \in E_u$ with $v \neq w$}
            \State Add edge $(v_{uv}, v_{uw})$ to $E_L^{(S)}$
        \EndFor
    \EndFor
    
    \State $L_S \gets \left( V_L^{(S)}, E_L^{(S)}, \mathbf{X}_L^{(S)}, \mathbf{Y}^{(S)} \right)$
\EndFor

\State \Return $(L_{\text{train}}, L_{\text{val}}, L_{\text{test}})$
\end{algorithmic}
\end{algorithm}


The resulting line graphs $L_S$ for each split $S \in \{\text{train}, \text{val}, \text{test}\}$. 
For each split $S$ is a tuple $L_S = (V_L^{(S)}, E_L^{(S)}, \mathbf{X}_L^{(S)}, \mathbf{Y}^{(S)})$, where $\mathbf{X}_L^{(S)}$ are the transformed node features and $\mathbf{Y}^{(S)}$ are the binary link labels, forming the direct input to the subsequent GNN model (Section~\ref{sec:gnn_architecture}).


\begin{lemma}[Complexity of Line Graph Construction]
\label{lemma:complexity}
Let $G=(V,E)$ be a graph with $n=|V|$ vertices and $m=|E|$ edges, where each vertex $v \in V$ has degree $d_v$, maximum degree $d_{\max}$, node features of dimension $d$, and optional edge features of dimension $d_e$. The line graph $L(G)$ constructed via Algorithm~\ref{algo:line_graph_generation} satisfies:

\begin{enumerate}
    \item \textbf{Time Complexity:} The construction requires $\Theta\left(\sum_{v \in V} \binom{d_v}{2}\right) = \Theta\left(\sum_{v \in V} d_v^2\right)$ operations for edge generation. Feature transformation (when applicable) adds $O(m \cdot (d + d_e))$ operations for computing the Hadamard product and optional edge gating. In the worst case, the total time is $O(m \cdot d_{\max} + m \cdot (d + d_e))$.
    
    \item \textbf{Space Complexity:} Storing $L(G)$ requires $\Theta\left(m + \sum_{v \in V} \binom{d_v}{2}\right)$ space for the graph structure. Feature storage (when applicable) requires $O(m \cdot d)$ space for the transformed node features, independent of edge features due to the gating mechanism's integration during transformation. The total space is $\Theta\left(m + \sum_{v \in V} d_v^2 + m \cdot d\right)$.
    
    \item \textbf{GNN Parameter Impact:} Unlike concatenation-based approaches, the Hadamard product maintains input dimension $d$, keeping the first-layer parameter count unchanged from processing original node features directly. When edge features are integrated, additional parameters are introduced only in the transformation layer ($\mathbf{W}_e \in \mathbb{R}^{d \times d_e}$, $\mathbf{b}_e \in \mathbb{R}^d$, and MLP parameters) rather than in the GNN's message-passing layers.
\end{enumerate}
\end{lemma}

\begin{proof}
\textbf{Time Complexity:} The algorithm processes each vertex $v \in V$ independently. For a vertex of degree $d_v$, it considers all unordered pairs of edges incident to $v$, of which there are $\binom{d_v}{2} = \frac{d_v(d_v-1)}{2}$. Summing over all vertices yields the total operation count for edge generation. Since $\binom{d_v}{2} \leq d_v^2$, the complexity is bounded by $\sum_{v \in V} d_v^2$. By noting that $\sum_{v \in V} d_v = 2m$ and applying the inequality $d_v \leq d_{\max}$, we obtain $\sum_{v \in V} d_v^2 \leq d_{\max} \sum_{v \in V} d_v = 2m \cdot d_{\max}$. Feature transformation requires processing each edge once: the Hadamard product takes $O(d)$ operations, and when edge features are available, the gating mechanism adds $O(d \cdot d_e)$ for the projection and $O(d_e)$ for the MLP, giving total $O(m \cdot (d + d_e))$.

\textbf{Space Complexity:} $L(G)$ contains exactly $m$ nodes (one per edge in $G$). Its edges correspond to shared endpoints in the original graph: for each vertex $v$, every pair of its incident edges becomes an edge in $L(G)$. Therefore, the total number of edges in $L(G)$ is $\sum_{v \in V} \binom{d_v}{2}$. Feature storage requires $d$ dimensions per node for $m$ nodes (regardless of edge feature integration), as the transformation outputs dimension $d$.

\textbf{GNN Parameters:} The Hadamard product transformation maintains the original feature dimension $d$, so the first linear layer in the GNN transforms features from dimension $d$ to the hidden dimension $h$, requiring a weight matrix of size $d \times h$, identical to processing original node features. Additional parameters from edge feature integration ($\mathbf{W}_e$, $\mathbf{b}_e$, MLP) are confined to the preprocessing transformation and do not affect GNN layer sizes.
\end{proof}

This approach maintains efficiency for sparse graphs where $d_{\max}$ is bounded. The Hadamard product's preservation of input dimension $d$ avoids parameter inflation in downstream GNN layers while capturing rich multiplicative interactions, making it particularly suitable for large-scale graphs with high-dimensional features.


\subsubsection{Theoretical Analysis of Line Graph Properties}

The line graph transformation provides several theoretical advantages that address key challenges in link prediction. Below we analyze three critical aspects: expressivity via the Weisfeiler-Lehman test, robustness to class imbalance, and adaptation to graph sparsity.
The expressive power of GNNs on line graphs can be formally characterized through the lens of the Weisfeiler-Lehman (WL) test, a standard measure for graph isomorphism.
The following theorem establishes the connection between edge discrimination in the original graph and node discrimination in its line graph.

\begin{theorem}[Expressivity of Line Graph Transformation]
\label{thm:wl_expressivity}
Let $G = (V, E)$ and $G' = (V', E')$ be two undirected graphs. For any edges $e = (u, v) \in E$ and $e' = (u', v') \in E'$, let $c^t_G(e)$ and $c^t_{G'}(e')$ denote the colors assigned to $e$ and $e'$ at iteration $t$ of the 1-Weisfeiler-Lehman (1-WL) test on $G$ and $G'$, respectively, where edge colors are derived from endpoint node colors as $c^t_G(e) = \tau(c^t_G(u), c^t_G(v))$ with $\tau$ being an injective function that encodes an unordered pair. If there exists an iteration $t$ such that $c^t_G(e) \neq c^t_{G'}(e')$, then there exists a graph neural network $\mathcal{F}$ operating on the line graphs $L(G)$ and $L(G')$ such that $\mathcal{F}(v_e) \neq \mathcal{F}(v_{e'})$, where $v_e$ and $v_{e'}$ are the nodes in $L(G)$ and $L(G')$ corresponding to $e$ and $e'$.

Furthermore, for any $k \geq 0$, the representation of $v_e$ computed by a $k$-layer GNN on $L(G)$ is determined by $c^k_G(e)$, i.e., the edge color at iteration $k$ of the 1-WL test on $G$. Consequently, a $k$-layer GNN on $L(G)$ is at least as expressive as $k$ iterations of the 1-WL test on $G$ for edge discrimination.
\end{theorem}

\begin{proof}
We prove the theorem in three parts.

\noindent\textbf{Part 1: Equivalence between 1-WL on $G$ and 1-WL on $L(G)$.}

Let $c^t_G(u)$ denote the color of node $u \in V$ at iteration $t$ of the 1-WL test on $G$, with $c^0_G(u)$ being the initial color (e.g., based on node degree). For each edge $e = (u, v) \in E$, define its color at iteration $t$ as:
\begin{equation}
    c^t_G(e) = \tau\left(c^t_G(u), c^t_G(v)\right),
\end{equation}


where $\tau$ is an injective function that encodes an unordered pair (e.g., a sorted tuple).

Now, consider the line graph $L(G) = (V_L, E_L)$ with $V_L = E$. We run the 1-WL test on $L(G)$ with initial colors for each node $v_e \in V_L$ set to $c^0_G(e)$. Let $d^t_{L(G)}(v_e)$ denote the color of $v_e$ at iteration $t$ of the 1-WL test on $L(G)$. We claim that for all $t \geq 0$:
\begin{equation}
d^t_{L(G)}(v_e) = c^t_G(e). 
\label{equ:weiProof1}
\end{equation}

    

We prove Equation~\ref{equ:weiProof1} by induction on $t$.

\textit{Base case ($t=0$):} By construction, $d^0_{L(G)}(v_e) = c^0_G(e)$.

\textit{Inductive step:} Assume Equation~\ref{equ:weiProof1} holds for some $t \geq 0$. Consider an edge $e = (u, v) \in E$. The neighbors of $v_e$ in $L(G)$ are exactly those edges in $G$ that share a common endpoint with $e$, i.e.,
\[
N_{L(G)}(v_e) = \{ v_{e'} \mid e' = (u, w) \in E, w \neq v \} \cup \{ v_{e'} \mid e' = (v, w) \in E, w \neq u \}.
\]
By the induction hypothesis, for each such neighbor $v_{e'}$, we have $d^t_{L(G)}(v_{e'}) = c^t_G(e')$. The 1-WL update on $L(G)$ is:
\[
d^{t+1}_{L(G)}(v_e) = \hash\left( d^t_{L(G)}(v_e), \{\{ d^t_{L(G)}(v_{e'}) : v_{e'} \in N_{L(G)}(v_e) \}\} \right),
\]
where $\hash$ is an injective hash function. Substituting the inductive hypothesis:
\begin{align*}
d^{t+1}_{L(G)}(v_e) &= \hash\left( c^t_G(e), \{\{ c^t_G(e') : e' \in N(e) \}\} \right) \\
&= \hash\Big( \tau(c^t_G(u), c^t_G(v)), \{\{ \tau(c^t_G(u), c^t_G(w)) : w \in N_G(u)\setminus\{v\} \}\} \\
&\quad \cup \{\{ \tau(c^t_G(v), c^t_G(w)) : w \in N_G(v)\setminus\{u\} \}\} \Big).
\end{align*}
Now, observe that the multiset $\{\{ c^t_G(e') : e' \in N(e) \}\}$ can be partitioned into two subsets: one containing pairs with first component $c^t_G(u)$ (edges incident to $u$) and another with first component $c^t_G(v)$ (edges incident to $v$). Since $\tau$ is injective, we can uniquely recover the multisets:
\[
M_u = \{\{ c^t_G(w) : w \in N_G(u) \}\} \quad \text{and} \quad M_v = \{\{ c^t_G(w) : w \in N_G(v) \}\}.
\]
Therefore,
\[
d^{t+1}_{L(G)}(v_e) = \hash\Big( \tau\big(c^t_G(u), c^t_G(v)\big), \tau\big(c^t_G(u), M_u\big), \tau\big(c^t_G(v), M_v\big) \Big).
\]
On the other hand, the 1-WL update on $G$ gives:
\[
c^{t+1}_G(u) = \hash(c^t_G(u), M_u), \quad c^{t+1}_G(v) = \hash(c^t_G(v), M_v).
\]
Thus,
\[
c^{t+1}_G(e) = \tau\left(c^{t+1}_G(u), c^{t+1}_G(v)\right) = \tau\left( \hash(c^t_G(u), M_u), \hash(c^t_G(v), M_v) \right).
\]
Since $\hash$ and $\tau$ are injective, the composition is injective, and we have:
\[
d^{t+1}_{L(G)}(v_e) = c^{t+1}_G(e).
\]
This completes the induction, proving Equation~\ref{equ:weiProof1}.

\noindent\textbf{Part 2: Existence of a distinguishing GNN.}

If $c^t_G(e) \neq c^t_{G'}(e')$ for some $t$, then by Equation~\ref{equ:weiProof1} $d^t_{L(G)}(v_e) \neq d^t_{L(G')}(v_{e'})$. It is known from \cite{xu2018how} that for any two graphs, if the 1-WL test assigns different colors to two nodes, then there exists a GNN that maps them to different representations. Therefore, there exists a GNN $\mathcal{F}$ such that $\mathcal{F}(v_e) \neq \mathcal{F}(v_{e'})$.

\noindent\textbf{Part 3: Expressivity of a $k$-layer GNN.}

A $k$-layer GNN on $L(G)$ computes a representation for $v_e$ that depends on the $k$-hop neighborhood of $v_e$ in $L(G)$. By Equation~\ref{equ:weiProof1}, the color $d^k_{L(G)}(v_e) = c^k_G(e)$ is determined by the $k$-hop neighborhood of $v_e$. Since the GNN can approximate the 1-WL test arbitrarily closely \cite{xu2018how}, the representation of $v_e$ is determined by $c^k_G(e)$. Hence, the GNN can distinguish edges $e$ and $e'$ whenever the 1-WL test on $G$ distinguishes them within $k$ iterations. Therefore, a $k$-layer GNN on $L(G)$ is at least as expressive as $k$ iterations of the 1-WL test on $G$ for edge discrimination.
\end{proof}

\begin{remark}
The theorem establishes a formal connection between the expressivity of GNNs on line graphs and the Weisfeiler-Lehman test on the original graph. This provides a theoretical foundation for the superior performance of LineML in capturing complex structural patterns. Moreover, it explains why the line graph transformation enhances discriminative power: by converting edges to nodes, we enable GNNs to leverage the full expressivity of the 1-WL test on the line graph, which is equivalent to a higher-order WL test on the original graph for edge discrimination.
\end{remark}

This theorem provides theoretical justification for using GNNs on line graphs: they inherit the strong discriminative power of the WL test while operating on a transformed graph where edges become first-class citizens.This result extends known equivalences between line graph WL tests and edge-WL \cite{bodnar2021weisfeiler,chen2021link} and the connection between GNNs and WL tests \cite{xu2018how}.



    
\subsection{Line Graph Pruning}
\label{subsec:linegraph_pruning}

The computational complexity of processing line graphs grows superlinearly with graph size, creating significant bottlenecks for large-scale applications. 
Thus, we propose a principled framework for line graph pruning that systematically reduces the number of edges while preserving essential structural and spectral properties of the graph. 
Pruning serves as both a compression technique and a regularization mechanism.
This can improve generalization by removing noisy or redundant connections. 



The pruning operation $\mathcal{F}: E_L \to E_L'$ maps the original edge set to a pruned subset $E_L' \subset E_L$ such that $|E_L'| = \eta |E_L|$, where $\eta \in (0,1]$ is the preservation ratio.

\begin{definition}[Pruning Quality Metric]
The quality of pruning is measured by the spectral distortion factor $\kappa$:
\begin{equation}
 \kappa(L(G), L(G)') = \max_{x \in \mathbb{R}^{|V_L|} \setminus \{0\}} \frac{x^T \mathbb{L} x}{x^T L' x} \cdot \max_{x \in \mathbb{R}^{|V_L|} \setminus \{0\}} \frac{x^T L' x}{x^T \mathbb{L} x}
\end{equation}
where $\mathbb{L}$ and $\mathbb{L}'$ are the Laplacians of $L(G)$ and $L(G)'$, respectively. We say pruning is $\epsilon$-spectrally preserving if $\kappa \leq 1 + \epsilon$.
\end{definition}

The spectral distortion factor $\kappa$ quantifies how much the Laplacian spectrum, and thus global properties such as connectivity and diffusion characteristics is altered by pruning. 
A value close to 1 indicates that the pruned graph is spectrally similar to the original, preserving critical structural invariants.

\subsubsection{k-Nearest Neighbor (kNN) Spectral Pruning}

 
The kNN pruning strategy is motivated by the observation that in feature-rich line graphs, edges connecting nodes with dissimilar feature vectors contribute less to the overall structural cohesion. 
By retaining only the top-$k$ most similar neighbors for each node, we preserve local neighborhood structures that are most informative for downstream tasks like link prediction. 
This approach is particularly effective when node embeddings contain meaningful semantic information that correlates with graph connectivity.


For each node in the line graph, we compute the Euclidean distance between its feature vector and those of its neighbors. 
For graphs lacking node attributes, these feature vectors are first generated using structural embeddings derived from methods like random walk matrix factorization.
We then retain only the $k$ nearest neighbors and remove all other edges. 
The process is symmetric to ensure undirectedness. 
This creates a sparser graph where connections are based on feature similarity rather than just structural adjacency.

Computing pairwise distances for each node requires $O(|V_L| \cdot \bar{d} \cdot f)$ time, where $\bar{d}$ is the average degree and $f$ is the feature dimension. Sorting neighbors takes $O(\bar{d} \log \bar{d})$ per node. The overall complexity is $O(|V_L| \cdot \bar{d} \cdot (f + \log \bar{d}))$, which is linear in the number of nodes and practically efficient for large graphs. Memory requirements are $O(|V_L| \cdot f)$ for storing feature vectors.

\begin{algorithm}[H]
\caption{kNN Spectral Pruning for Line Graphs}
\label{algo:knn_pruning}
\begin{algorithmic}[1]
\Require Line graph $L(G) = (V_L, E_L)$ with node features $X \in \mathbb{R}^{|V_L| \times d}$, parameter $k$
\Ensure Pruned line graph $L(G)' = (V_L, E_L')$
\State Initialize $E_L' \gets \emptyset$
\For{each node $v \in V_L$}
    \State Let $N(v) = \{u : (u,v) \in E_L\}$ be neighbor set
    \State Compute feature distance: $d_{uv} = \|X_u - X_v\|_2$ for $u \in N(v)$
    \State Sort neighbors by distance: $N_{\text{sorted}}(v) \gets \text{sort}(N(v), \text{key}=d_{uv})$
    \State Select top-$k$ nearest neighbors: $N_k(v) \gets N_{\text{sorted}}(v)[:k]$
    \State Add edges to $E_L'$: $E_L' \gets E_L' \cup \{(v,u) : u \in N_k(v)\}$
\EndFor
\State Symmetrize: $E_L' \gets \{(u,v) : (u,v) \in E_L' \text{ or } (v,u) \in E_L'\}$
\State \Return $L(G)' = (V_L, E_L')$
\end{algorithmic}
\end{algorithm}

\begin{theorem}[kNN Pruning Spectral Perturbation]
For an unweighted line graph $L(G)$ with maximum degree $\Delta$, let $L(G)'$ be the graph after kNN pruning with parameter $k$. Let:
- $d_{\text{removed}}(u) = |\{v: (u,v) \in E_{\text{removed}}\}|$ be the number of edges removed incident to node $u$
- $\lambda_2(L(G))$ be the second smallest eigenvalue (algebraic connectivity) of $L(G)$

If $\max_u d_{\text{removed}}(u) < \lambda_2(L(G))$, then:
\[
\kappa(L(G), L(G)') \leq 1 + \frac{\max_u d_{\text{removed}}(u)}{\lambda_2(L(G)) - \max_u d_{\text{removed}}(u)}.
\]
Moreover, since $d_{\text{removed}}(u) \leq \Delta - k$ for all $u$, we have:
\[
\kappa(L(G), L(G)') \leq 1 + \frac{\Delta - k}{\lambda_2(L(G)) - (\Delta - k)} \quad \text{when } \Delta - k < \lambda_2(L(G)).
\]
\end{theorem}

\begin{proof}
Let $L$ and $L'$ be the Laplacians of $L(G)$ and $L(G)'$ respectively. Write $L = L' + R$, where $R$ is the Laplacian of the graph consisting only of removed edges.

\textbf{Step 1: Relating the condition number to the perturbation.}
For any nonzero $x \in \mathbb{R}^{|V_L|}$, we have:
\[
\frac{x^T L x}{x^T L' x} = 1 + \frac{x^T R x}{x^T L' x}.
\]
Let $\mu_{\max} = \max_{x \neq 0} \frac{x^T L x}{x^T L' x}$ and $\mu_{\min} = \min_{x \neq 0} \frac{x^T L x}{x^T L' x}$. Then $\kappa = \mu_{\max} / \mu_{\min}$.

\textbf{Step 2: Bounding $\mu_{\max}$.}
Since $R$ is positive semidefinite, $x^T R x \geq 0$ for all $x$. Thus:
\[
\mu_{\max} = 1 + \max_{x \neq 0} \frac{x^T R x}{x^T L' x}.
\]
By the Courant-Fischer theorem for the generalized Rayleigh quotient:
\[
\max_{x \neq 0} \frac{x^T R x}{x^T L' x} = \lambda_{\max}(L'^{-1/2} R L'^{-1/2}) \leq \frac{\|R\|_2}{\lambda_{\min}(L')},
\]
where $\|R\|_2$ is the spectral norm of $R$ and $\lambda_{\min}(L')$ is the smallest positive eigenvalue of $L'$.

\textbf{Step 3: Bounding $\|R\|_2$ and $\lambda_{\min}(L')$.}
For the Laplacian $R$ of a graph where each node has at most $d_{\max} = \max_u d_{\text{removed}}(u)$ incident edges, we have $\|R\|_2 \leq 2d_{\max} \leq 2(\Delta - k)$ by Gershgorin circle theorem.

By Weyl's inequality for Laplacians, $\lambda_{\min}(L') \geq \lambda_2(L) - \|R\|_2 \geq \lambda_2(L) - d_{\max}$.

\textbf{Step 4: Bounding $\mu_{\min}$.}
Similarly,
\[
\mu_{\min} = 1 + \min_{x \neq 0} \frac{x^T R x}{x^T L' x} \geq 1,
\]
since $x^T R x \geq 0$.

\textbf{Step 5: Combining bounds.}
Thus,
\[
\kappa = \frac{\mu_{\max}}{\mu_{\min}} \leq \frac{1 + \frac{d_{\max}}{\lambda_2(L) - d_{\max}}}{1} = 1 + \frac{d_{\max}}{\lambda_2(L) - d_{\max}}.
\]
Substituting $d_{\max} \leq \Delta - k$ gives the second inequality.

\textbf{Step 6: Tightening the bound (optional).}
A tighter bound can be obtained using the fact that for Laplacians, $\|R\|_2 \leq \max_u d_{\text{removed}}(u) + \max_v d_{\text{removed}}(v) = 2d_{\max}$, but the actual spectral norm is often closer to $d_{\max}$. The simplified bound uses $\|R\|_2 \leq d_{\max}$ which holds since the maximum eigenvalue of a Laplacian is at most the maximum degree.
\end{proof}

\begin{remark}
The condition $\Delta - k < \lambda_2(L(G))$ ensures the bound is meaningful. For well-connected graphs where $\lambda_2(L(G))$ is large relative to $\Delta$, kNN pruning can remove significant numbers of edges while maintaining spectral similarity. In practice, even when this condition doesn't hold exactly, the theorem provides qualitative insight: pruning preserves spectral structure better when:
\begin{enumerate}
    \item The original graph has high algebraic connectivity (well-connected)
    \item $k$ is not too small relative to $\Delta$ (preserving enough edges per node)
    \item The removed edges are distributed evenly across nodes (minimizing $\max_u d_{\text{removed}}(u)$)
\end{enumerate}
\end{remark}

\subsubsection{Spectral Sparsification via Effective Resistance Sampling}

Spectral sparsification aims to produce a subgraph that preserves the Laplacian quadratic form $x^\top \mathbf{M} x$.
This method governs essential graph properties including connectivity, mixing time, and spectral clustering \cite{spielman2011graph}.
The core idea is to select a subset of edges whose weighted Laplacian approximates the original. 
Effective resistance $\Omega_{uv}$, defined as the electrical resistance between nodes $u$ and $v$ when the graph is viewed as a resistor network, quantifies the importance of each edge in maintaining global connectivity. 
Edges with high effective resistance act as bridges connecting disparate components and are therefore critical for preserving spectral properties \cite{spielman2011graph}.



Exact effective resistance computation requires $O(|V_L|^3)$ time via matrix inversion, which is prohibitive for large graphs. We therefore employ the nearly-linear time approximation of Spielman and Srivastava \cite{spielman2011graph}, which uses Johnson-Lindenstrauss random projections to estimate all pairwise effective resistances simultaneously. This reduces complexity to $O(|E_L| \cdot r)$ where $r = O(\log |V_L|/\epsilon^2)$ is the projection dimension, with memory requirements of $O(|V_L| \cdot r)$ for the projection matrix.

\begin{definition}[Effective Resistance Approximation]
Given line graph $L(G) = (V_L, E_L)$ with Laplacian matrix $\mathbf{M} \in \mathbb{R}^{|V_L| \times |V_L|}$, the effective resistance between nodes $u$ and $v$ is:
\[
\Omega_{uv} = (e_u - e_v)^\top \mathbf{M}^\dagger (e_u - e_v),
\]
where $\mathbf{M}^\dagger$ denotes the Moore-Penrose pseudoinverse and $e_u$ is the standard basis vector. Following \cite{spielman2011graph, cohen2015nearly}, we approximate $\Omega_{uv}$ via Johnson-Lindenstrauss projection:
\[
\tilde{\Omega}_{uv} = \|\mathbf{Q}(e_u - e_v)\|_2^2,
\]
where $\mathbf{Q} \in \mathbb{R}^{r \times |V_L|}$ is a random matrix with entries $\pm 1/\sqrt{r}$.
\end{definition}


The parameter $\eta \in (0, 1]$ controls the sparsity of the output graph, representing the target fraction of edges to retain. 
In Algorithm \ref{algo:spectral_sparsification}, $\eta$ influences the sampling process: 
edges are sampled with probability $p_{uv} \propto \tilde{\Omega}_{uv} \cdot w_{uv}$, 
then accepted with probability $\min(1, C \cdot p_{uv} \cdot \log |V_L| / \epsilon^2)$. 
The constant $C$ is adjusted iteratively to achieve the target edge count $|E_L'| \approx \eta |E_L|$. 
This provides explicit control over the trade-off between computational efficiency (lower $\eta$) and spectral fidelity (higher $\eta$).

\begin{algorithm}[!ht]
\caption{Approximate Spectral Sparsification with Preservation Control}
\label{algo:spectral_sparsification}
\begin{algorithmic}[1]
\Require Line graph $L(G) = (V_L, E_L)$ with edge weights $\{w_{uv}\}$
\Require Preservation ratio $\eta \in (0, 1]$, tolerance $\epsilon > 0$, failure probability $\delta > 0$
\Ensure $\epsilon$-spectrally similar sparsifier $L(G)'$ with $|E_L'| \approx \eta |E_L|$
\State Set projection dimension $r \gets \lceil 4\log(|V_L|/\delta)/\epsilon^2 \rceil$
\State Generate random projection matrix $\mathbf{Q} \in \mathbb{R}^{r \times |V_L|}$ with $\mathbf{Q}_{ij} \sim \{\pm 1/\sqrt{r}\}$
\For{each edge $(u,v) \in E_L$}
    \State Compute projection: $\tilde{\Omega}_{uv} \gets \|\mathbf{Q}(e_u - e_v)\|_2^2$ \Comment{Approx. effective resistance}
    \State Set sampling probability: $p_{uv} \gets \frac{\tilde{\Omega}_{uv} \cdot w_{uv}}{\sum_{(x,y) \in E_L} \tilde{\Omega}_{xy} \cdot w_{xy}}$
\EndFor
\State Initialize scale factor $C \gets 1.0$
\Repeat
    \For{each edge $(u,v) \in E_L$}
        \State Compute acceptance probability: $q_{uv} \gets \min(1, C \cdot p_{uv} \cdot \log |V_L| / \epsilon^2)$
        \State Sample edge with probability $q_{uv}$, add to $E_L'$ if selected
    \EndFor
    \State Update $C \gets C \cdot \frac{\eta |E_L|}{|E_L'|}$ \Comment{Adjust to hit target $\eta$}
\Until{$| |E_L'| - \eta |E_L| | < 0.01 \eta |E_L|$} \Comment{Convergence condition}
\For{each selected edge $(u,v) \in E_L'$}
    \State Set new weight: $w_{uv}' \gets w_{uv} / q_{uv}$ \Comment{Compensatory reweighting}
\EndFor
\State \Return $L(G)' = (V_L, E_L', \{w_{uv}'\})$
\end{algorithmic}
\end{algorithm}


Algorithm \ref{algo:spectral_sparsification} provides the following guarantee: with probability at least $1-\delta$, the output graph $L(G)'$ satisfies the spectral approximation condition
\[
(1 - \epsilon) x^\top \mathbf{M} x \leq x^\top \mathbf{M}' x \leq (1 + \epsilon) x^\top \mathbf{M} x \quad \forall x \in \mathbb{R}^{|V_L|},
\]
where $\mathbf{M}$ and $\mathbf{M}'$ are the Laplacians of $L(G)$ and $L(G)'$, respectively. This guarantee ensures that key spectral properties are preserved within multiplicative factor $(1\pm\epsilon)$. The bound follows directly from the seminal work on graph sparsification via effective resistance sampling \cite{spielman2011graph}, with the Johnson-Lindenstrauss approximation preserving effective resistances within factor $(1\pm\epsilon/2)$ with high probability \cite{cohen2015nearly}.


\subsubsection{Degree-Based Hierarchical Pruning}


Many real-world networks exhibit power-law degree distributions, where a small fraction of nodes (hubs) have very high degree. 
Degree-based pruning leverages this heterogeneity: removing edges incident to high-degree nodes can drastically reduce edge density while preserving connectivity due to the redundant connections of hubs. 
Conversely, low-degree nodes often contribute less to global connectivity and can be pruned with minimal impact. 
This approach directly addresses memory bandwidth limitations in HPC systems by reducing adjacency matrix density.

We offer two complementary strategies:  
\textbf{High-degree pruning}: Remove edges from nodes whose normalized degree exceeds a threshold $\tau$. This targets hubs that contribute disproportionately to edge count.  
\textbf{Low-degree pruning}: Remove edges from nodes with normalized degree below $\tau$. This eliminates peripheral connections with minimal structural impact.  
A BFS spanning tree is then added to ensure the pruned graph remains connected, guaranteeing at least one path between any two nodes.


Degree computation is $O(|V_L|)$. Edge removal is $O(|E_L|)$. The BFS step is $O(|V_L| + |E_L'|)$. 
Overall complexity is linear in the size of the line graph. Memory requirements are minimal, requiring only degree storage and edge removal tracking.

\begin{definition}[Degree Concentration]
For line graph $L(G)$, let $D = \text{diag}(d_1, \ldots, d_{|V_L|})$ be the degree matrix. The degree concentration ratio is:
\[
\rho = \frac{\max_i d_i}{\min_i d_i}
\]
High $\rho$ indicates heterogeneous degree distribution amenable to degree-based pruning.
\end{definition}

\begin{algorithm}[H]
\caption{Degree-Based Hierarchical Pruning}
\label{algo:degree_pruning}
\begin{algorithmic}[1]
\Require Line graph $L(G)$, threshold $\tau$, mode $\in \{\text{high}, \text{low}\}$
\Ensure Pruned line graph $L(G)'$
\State Compute normalized degrees: $\tilde{d}_i = d_i / \max_j d_j$
\If{mode = 'high'}
    \State Remove edges incident to nodes with $\tilde{d}_i > \tau$
\ElsIf{mode = 'low'}
    \State Remove edges incident to nodes with $\tilde{d}_i < \tau$
\EndIf
\State Ensure connectivity via BFS spanning tree
\State \Return $L(G)'$
\end{algorithmic}
\end{algorithm}


Pruning reduces both computational and communication overhead, especially in distributed HPC settings. The following analysis quantifies the training time reduction.
\begin{lemma}[Training Time Reduction]
\label{lemma:trainingTimeReduction}
Let $T_{\text{orig}}$ be the training time per epoch on the original line graph $L(G)$. After pruning with preservation ratio $\eta$, the training time per epoch reduces to:
\[
T_{\text{pruned}} = O(\eta \cdot T_{\text{orig}}).
\]
Furthermore, in distributed settings with $P$ processors, the communication overhead reduces proportionally to $\eta$, and the improved cache utilization can yield additional practical speedups.
\end{lemma}

\begin{proof}
The computational complexity of a GNN forward pass scales linearly with the number of edges processed. For a $K$-layer GNN with feature dimension $d$, the per-epoch training time satisfies $T_{\text{orig}} = \Theta(|E_L| \cdot K \cdot d)$. Pruning reduces the edge set from $|E_L|$ to approximately $\eta|E_L|$, yielding $T_{\text{comp}}^{\text{pruned}} = \Theta(\eta|E_L| \cdot K \cdot d) = \eta \cdot T_{\text{orig}}$.

In distributed training, communication overhead typically scales with the number of edges crossing partition boundaries, denoted $|E_{\text{cut}}|$. Under uniform pruning, the number of cut edges reduces to approximately $\eta|E_{\text{cut}}|$, decreasing communication time proportionally. Additionally, the reduced graph size improves cache locality, often yielding faster memory access in practice. Combining these effects gives the overall $O(\eta \cdot T_{\text{orig}})$ scaling.
\end{proof}





\subsection{Graph Neural Network Architecture}
\label{sec:gnn_architecture}

This section details the Graph Neural Network architecture that forms the core of our model. 
We adopt the GraphSAGE framework \cite{hamilton2017inductive} for its inductive capability to generate node embeddings by learning aggregation functions from a node's local neighborhood. 
This design enables the model to synthesize both the contextual attributes of nodes and the geometric connectivity of the graph structure, producing expressive embeddings ideal for downstream link prediction.

The model processes the constructed line graph $L(G) = (V_L, E_L)$. 
Let $\mathbf{X}_L \in \mathbb{R}^{|V_L| \times d_{in}}$ denote the initial feature matrix for all nodes in $V_L$, where $d_{in}$ is the input feature dimension. 
The architecture is composed of $K$ successive GraphSAGE layers. 
Each layer $k$ transforms the node representations by aggregating information from their topological neighbors.


At each layer $k$, the model updates the representation $\mathbf{h}_v^{(k)}$ for a node $v \in V_L$. The update rule is defined as:
\begin{equation}
    \mathbf{h}_v^{(k)} = \sigma \left( \mathbf{W}^{(k)} \cdot \left[ \mathbf{h}_v^{(k-1)} \parallel \text{AGGREGATE}^{(k)} \left( \{ \mathbf{h}_u^{(k-1)} : u \in \mathcal{N}(v) \} \right) \right] \right),
\end{equation}
where:
\begin{itemize}
    \item $\mathbf{h}_v^{(k-1)}$ is the representation of node $v$ from the previous layer $(k-1)$, with $\mathbf{h}_v^{(0)} = \mathbf{x}_v$ being its initial feature vector.
    \item $\mathcal{N}(v)$ is the set of neighboring nodes of $v$ in $L(G)$.
    \item $\text{AGGREGATE}^{(k)}(\cdot)$ is a differentiable, permutation-invariant function that consolidates features from the neighbor set.
    \item $\mathbf{W}^{(k)}$ is a layer-specific trainable weight matrix.
    \item $\sigma(\cdot)$ denotes a non-linear activation function (ReLU).
    \item $\parallel$ denotes the vector concatenation operation.
\end{itemize}

The aggregation function is crucial for capturing neighborhood structures. GraphSAGE supports several variants:
\begin{itemize}
    \item \textbf{Mean Aggregation:} Computes the element-wise mean of neighbor embeddings: $\text{AGGREGATE}_{\text{mean}} = \frac{1}{|\mathcal{N}(v)|} \sum_{u \in \mathcal{N}(v)} \mathbf{h}_u^{(k-1)}$.
    \item \textbf{Max-Pooling Aggregation:} Applies a learned transformation followed by element-wise max pooling: $\text{AGGREGATE}_{\text{pool}} = \max( \{\sigma(\mathbf{W}_{\text{pool}} \mathbf{h}_u^{(k-1)} + \mathbf{b}) \mid u \in \mathcal{N}(v) \} )$.
    \item \textbf{LSTM Aggregation:} Uses a Long Short-Term Memory network to process a random permutation of the neighbors: $\text{AGGREGATE}_{\text{LSTM}} = \text{LSTM}([ \mathbf{h}_u^{(k-1)}, \forall u \in \pi(\mathcal{N}(v)) ])$.
\end{itemize}


A stack of $K$ such layers allows information to propagate across $K$-hop neighborhoods in the line graph. The representation of a node after the final layer, $\mathbf{z}_v = \mathbf{h}_v^{(K)} \in \mathbb{R}^{d_{out}}$, encapsulates both its intrinsic features and the multi-scale structural context of its surroundings. To promote training stability and representational consistency, these final embeddings are  $\ell_2$-normalized:
\begin{equation}
    \mathbf{z}_v \leftarrow \frac{\mathbf{z}_v}{\|\mathbf{z}_v\|_2}.
\end{equation}

This sequence of transformations maps the initial features $\mathbf{X}_L$ of the line graph nodes to a set of refined, structure-aware embeddings $\mathbf{Z} = \{ \mathbf{z}_v : v \in V_L \}$. 
These embeddings $\mathbf{Z}$ serve as the direct input to the subsequent metric learning objective.

%
\subsection{Adaptive Triplet Margin for Informed Metric Learning}
\label{subsec:metric_learning}

Effective link prediction requires an embedding space where the representations of connected nodes are closer than those of unconnected nodes. 
To this end, we employ a triplet loss framework, a cornerstone of metric learning \cite{boutros2021self}. 
The learning signal is derived from triplets $(a, p, n)$, where $a$ is an anchor node in the line graph, $p$ is a positive node, and $n$ is a negative node.
The model is trained to pull $a$ and $p$ together while pushing $a$ and $n$ apart.

Let $f: V_L \to \mathbb{R}^d$ be the embedding function producing normalized node representations $\mathbf{z} = f(\cdot)$, and let $d(\mathbf{z}_i, \mathbf{z}_j) = \|\mathbf{z}_i - \mathbf{z}_j\|_2$ denote the Euclidean distance. The standard triplet loss enforces a fixed margin $\alpha$ between positive and negative distances:
\begin{equation}
\label{eq:standard_triplet}
    \mathcal{L}_{\text{std}} = \sum_{(a, p, n) \in \mathcal{T}} \left[ d(\mathbf{z}_a, \mathbf{z}_p) - d(\mathbf{z}_a, \mathbf{z}_n) + \alpha \right]_+,
\end{equation}
where $[x]_+ = \max(x, 0)$ and $\mathcal{T}$ is a set of sampled triplets.

A critical limitation of this formulation is its uniformity: it applies the same force of separation $\alpha$ to all triplets, regardless of their inherent difficulty. In graph learning, triplets exhibit a wide heterogeneity in difficulty; some negative edges are topologically or attribute-wise very similar to positive edges, making them ``hard" negatives that require a stronger push to separate.

To address this, we introduce a simple yet effective adaptive margin. Our key insight is to modulate the learning signal based on the triplet's current configuration, effectively applying a stronger gradient signal to harder examples. We define the \textit{adaptive triplet loss} as:
\begin{equation}
\label{eq:adaptive_triplet}
    \mathcal{L}_{\text{adapt}} = \sum_{(a, p, n) \in \mathcal{T}} \left[ (1 + \beta) \big( d(\mathbf{z}_a, \mathbf{z}_p) - d(\mathbf{z}_a, \mathbf{z}_n) \big) + \alpha \right]_+,
\end{equation}
where $\beta \ge 0$ is a global \textit{adaptivity hyperparameter}. This formulation maintains the fixed baseline margin $\alpha$ but scales the relative distance term by $(1+\beta)$. The scaling directly amplifies the gradient for active triplets, as established in the following proposition.

\begin{proposition}[Gradient Scaling and Active Set]
\label{prop:gradient_properties}
Let $\mathcal{A}_{\text{std}} = \{(a,p,n) : d_{ap} - d_{an} + \alpha > 0\}$ be the set of ``active" triplets for the standard loss where the gradient is non-zero. For the adaptive loss, the active set is $\mathcal{A}_{\text{adapt}} = \{(a,p,n) : (1+\beta)(d_{ap} - d_{an}) + \alpha > 0\} = \{(a,p,n) : d_{ap} - d_{an} + \alpha/(1+\beta) > 0\}$. 

For any triplet in $\mathcal{A}_{\text{adapt}}$, the gradient of the adaptive loss with respect to the distance difference $\delta = d_{ap} - d_{an}$ is scaled by a factor of $(1+\beta)$ compared to the standard loss:
\begin{equation}
    \frac{\partial \mathcal{L}_{\text{adapt}}}{\partial \delta} = (1+\beta) \cdot \frac{\partial \mathcal{L}_{\text{std}}}{\partial \delta}.
\end{equation}

Furthermore, $\mathcal{A}_{\text{adapt}} \subseteq \mathcal{A}_{\text{std}}$, with strict inclusion when $\beta > 0$ (i.e., $\mathcal{A}_{\text{adapt}} \subset \mathcal{A}_{\text{std}}$). This means the adaptive loss focuses on a subset of harder triplets while amplifying their gradients.
\end{proposition}

\begin{proof}
The gradient of the hinge-style loss $[x]_+$ with respect to its argument is $\mathbb{I}(x > 0)$, where $\mathbb{I}$ is the indicator function. For the standard loss, the argument is $\delta + \alpha$. Its gradient w.r.t. $\delta$ is therefore $\mathbb{I}(\delta + \alpha > 0)$. For the adaptive loss, the argument is $(1+\beta)\delta + \alpha$. Its gradient w.r.t. $\delta$ is $(1+\beta) \cdot \mathbb{I}((1+\beta)\delta + \alpha > 0)$.

For any triplet where $(1+\beta)\delta + \alpha > 0$ (i.e., $\mathcal{A}_{\text{adapt}}$), it follows that $\delta + \alpha > 0$ since $(1+\beta)\delta + \alpha > 0 \Rightarrow \delta > -\alpha/(1+\beta) \geq -\alpha$ for $\beta \geq 0$. Thus, $\mathcal{A}_{\text{adapt}} \subseteq \mathcal{A}_{\text{std}}$. For triplets in $\mathcal{A}_{\text{adapt}}$, both indicator functions equal 1, giving the gradient ratio $(1+\beta)$.

The inclusion is strict for $\beta>0$ because triplets with $\delta \in (-\alpha, -\alpha/(1+\beta))$ are in $\mathcal{A}_{\text{std}}$ but not in $\mathcal{A}_{\text{adapt}}$.
\end{proof}

Proposition~\ref{prop:gradient_properties} clarifies the mechanism of adaptivity: 
(1) \textbf{Gradient Amplification:} For the subset of triplets that are active under both losses, the adaptive loss amplifies their gradient by $(1+\beta)$. 
(2) \textbf{Focused Active Set:} The adaptive loss restricts learning to a subset of harder triplets where $d_{ap} - d_{an} > -\alpha/(1+\beta)$, ignoring easier triplets that are in $\mathcal{A}_{\text{std}}$ but have $\delta \in (-\alpha, -\alpha/(1+\beta)]$. The hyperparameter $\beta$ thus controls a trade-off: increasing $\beta$ focuses learning on progressively harder examples while amplifying their gradients, potentially refining decision boundaries around challenging cases.

Given this theoretical framework, the efficacy of triplet learning depends on sampling informative triplets $\mathcal{T}$ that are active under the adaptive loss. We therefore employ a targeted online mining strategy focused on generating \textbf{semi-hard negatives}. For each anchor $a$ corresponding to a positive edge, we sample a negative $n$ such that the triplet $(a, p, n)$ satisfies $d(\mathbf{z}_a, \mathbf{z}_p) - d(\mathbf{z}_a, \mathbf{z}_n) > -\alpha/(1+\beta)$ and $d(\mathbf{z}_a, \mathbf{z}_n) > d(\mathbf{z}_a, \mathbf{z}_p)$. This condition ensures the negative is farther from the anchor than the positive, yet the separation does not satisfy the fixed margin $\alpha$ by a comfortable amount, placing $\delta = d_{ap} - d_{an}$ in the interval $(-\alpha/(1+\beta), 0)$. These triplets are guaranteed to be active under $\mathcal{L}_{\text{adapt}}$ and receive an amplified gradient of $(1+\beta)$. By focusing on negatives that are already farther than the positive (avoiding the hardest negatives initially), we promote more stable training and help prevent collapse of the embedding space. To ensure diversity and mitigate sampling bias, this semi-hard mining is complemented by random sampling from the set of nodes not connected to $a$, with a bias towards higher-degree nodes to approximate topologically challenging negatives. This hybrid strategy ensures a steady and pedagogically sound supply of triplets, allowing the model to progressively learn finer distinctions in the embedding space.

The optimization of $\mathcal{L}_{\text{adapt}}$ directly refines the initial embeddings. The result of this process is the final, metric-aware embedding matrix:
\begin{equation}
    \mathbf{Z}_{\text{final}} = \{ \mathbf{z}_{v, \text{final}} : v \in V_L \},
\end{equation}
where $\mathbf{z}_{v, \text{final}}$ is the refined representation of node $v$. $\mathbf{Z}_{\text{final}}$ forms the direct input to the link classifier (Section~\ref{sec:link_classifier}).

\subsection{Link Classification Model}
\label{sec:link_classifier}

The final stage of the LineML framework is a dedicated link classifier that maps the refined, metric-aware embeddings $\mathbf{Z}_{\text{final}}$ from the adaptive metric learning module (Section~\ref{subsec:metric_learning}) to a probability of link existence. This module operates on the enhanced representation $\mathbf{z}_{v,\text{final}} \in \mathbb{R}^{d_{\text{out}}}$ for each node $v$ in the line graph $L(G)$.

To mitigate variance and enhance generalization, we employ an ensemble of $M$ independent multilayer perceptron (MLP) heads. Each head is a feed-forward network with non-linear transformations. For the $m$-th head, the transformation for a node $v$ is:
\begin{equation}
    \mathbf{g}_v^{(m, 1)} = \text{ReLU}\left( \text{LayerNorm}\left( \mathbf{W}_1^{(m)} \mathbf{z}_{v,\text{final}} + \mathbf{b}_1^{(m)} \right) \right),
\end{equation}
\begin{equation}
    \mathbf{g}_v^{(m, 2)} = \text{ReLU}\left( \mathbf{W}_2^{(m)} \mathbf{g}_v^{(m, 1)} + \mathbf{b}_2^{(m)} \right),
\end{equation}
\begin{equation}
    \hat{y}_v^{(m)} = \sigma\left( \mathbf{W}_3^{(m)} \mathbf{g}_v^{(m, 2)} + \mathbf{b}_3^{(m)} \right),
\end{equation}
where $\mathbf{W}_*^{(m)}, \mathbf{b}_*^{(m)}$ are learnable parameters unique to head $m$, $\text{LayerNorm}$ denotes layer normalization for training stability, $\text{ReLU}$ is the rectified linear unit activation, and $\sigma$ is the sigmoid function outputting a probability in $[0,1]$. Dropout is applied after the first layer for regularization.

The final prediction $\hat{y}_v$ is obtained by averaging the outputs of all $M$ heads:
\begin{equation}
    \hat{y}_v = \frac{1}{M} \sum_{m=1}^{M} \hat{y}_v^{(m)}.
\end{equation}
This ensemble approach reduces prediction variance and the risk of overfitting to the idiosyncrasies learned by any single classifier head.

The classifier is trained by minimizing the Binary Cross-Entropy (BCE) loss. For a batch of $B$ samples with ground-truth labels $y_v \in \{0,1\}$, the classification loss is:
\begin{equation}
    \mathcal{L}_{\text{class}} = - \frac{1}{B} \sum_{v=1}^{B} \left[ y_v \log(\hat{y}_v) + (1 - y_v) \log(1 - \hat{y}_v) \right].
\end{equation}

The entire LineML architecture is trained jointly in an end-to-end manner.  This joint optimization ensures the model learns a unified representation that is simultaneously \textit{structurally informative} (via the GNN), \textit{geometrically discriminative} (via metric learning), and directly \textit{predictive} of link existence.


\section{Experimental Study}
\label{sec:experiments}


This section presents a comprehensive empirical evaluation of the proposed LineML framework.
Our experimental analysis is structured around two primary dimensions: \textbf{functional evaluation} measuring predictive performance and \textbf{non-functional evaluation} assessing computational efficiency and scalability.

All experiments are conducted on a dedicated high-performance computing system equipped with dual NVIDIA RTX A4000 GPUs with 16GB VRAM each, an Intel Xeon Gold 5418Y processor, with 24 cores and 48 threads, and 128GB of system memory. 
The software environment utilizes Python 3.11.14 with PyTorch 2.6.0 and PyTorch Geometric 2.7.0 frameworks, leveraging CUDA 12.4 and cuDNN 9.1.0 for GPU acceleration. 
The system runs Ubuntu 24.04 LTS with NVIDIA driver version 580.95.05. 
This configuration ensures that all performance comparisons reflect realistic, production-ready computational capabilities rather than idealized or constrained environments.


\subsection{Benchmark Datasets and Statistical Characterization}
\label{sec:5.1}

To rigorously evaluate the proposed \textit{LineML} framework, we conduct experiments on a collection of 18 networks spanning multiple domains, including citation graphs, co-authorship networks, biological interaction networks, social systems, and infrastructure graphs. 
Most of the datasets are obtained from the Noesis link prediction repository \cite{watts1998collective,newman2001structure}\footnote{\url{https://noesis.ikor.org/datasets/link-prediction}}, which provides a standardized benchmark suite widely used in network analysis and link prediction studies.
In addition, we include three widely adopted citation benchmarks, \textit{Cora}, \textit{CiteSeer}, and \textit{PubMed}, introduced by Sen et al.~\cite{sen2008collective}, together with the \textit{Physics} co-authorship network from the KDD Cup collection~\cite{gehrke2003overview}. 
Table~\ref{tab:dataset} summarizes detailed statistics of both the original graphs and their corresponding line graphs.




\begin{table*}[!ht]
\centering
\caption{Statistics of Original Graphs and their corresponding Line Graphs (LG).}
\label{tab:dataset}
\resizebox{\textwidth}{!}{%
\begin{tabular}{lllrrrrrrr}
\toprule
\multicolumn{3}{c}{\textbf{Dataset Info}} & \multicolumn{4}{c}{\textbf{Original Graph}} & \multicolumn{3}{c}{\textbf{Line Graph (LG)}} \\
\cmidrule(lr){1-3} \cmidrule(lr){4-7} \cmidrule(lr){8-10}
\textbf{Category} & \textbf{Name} & \textbf{Type} & \textbf{Nodes} & \textbf{Edges} & \textbf{Feat.} & \textbf{Dens/Spar(\%)} & \textbf{Nodes} & \textbf{Edges} & \textbf{Dens/Spar(\%)} \\
\midrule

% --- Citation Networks ---
\multirow{3}{*}{\textbf{Citation}} 
 & Cora & Undirected & 2,708 & 5,278 & 1,433 & 0.14 / 99.86 & 21,112 & 314,992 & 0.07 / 99.93 \\
 & CiteSeer & Undirected & 3,327 & 4,552 & 3,703 & 0.08 / 99.92 & 18,208 & 203,788 & 0.06 / 99.94 \\
 & PubMed & Undirected & 19,717 & 44,324 & 500 & 0.02 / 99.98 & 177,296 & 3,496,168 & 0.01 / 99.99 \\
\midrule

% --- Co-authorship Networks ---
\multirow{7}{*}{\textbf{Co-authorship}} 
 & Physics & Undirected & 34,493 & 247,962 & 8,415 & 0.04 / 99.96 & 991,848 & 45,856,498 & 0.005 / 99.99 \\
 & ZWL & Directed & 6,651 & 108,364 & 0 & 0.25 / 99.75 & 216,728 & 11,873,265 & 0.03 / 99.97 \\
 & LDG & Directed & 8,324 & 83,064 & 0 & 0.12 / 99.88 & 166,128 & 9,409,089 & 0.03 / 99.97 \\
 & GRQ & Directed & 5,241 & 28,968 & 0 & 0.11 / 99.89 & 57,936 & 1,231,571 & 0.04 / 99.96 \\
 & NSC & Directed & 1,461 & 5,484 & 0 & 0.26 / 99.74 & 9,445 & 106,668 & 0.12 / 99.88 \\
 & SMG & Directed & 1,024 & 9,832 & 0 & 0.94 / 99.06 & 19,664 & 787,630 & 0.20 / 99.80 \\
 & KHN & Directed & 3,772 & 25,436 & 0 & 0.18 / 99.82 & 50,872 & 2,799,022 & 0.11 / 99.89 \\
\midrule

% --- Biological Networks ---
\multirow{3}{*}{\textbf{Biological}} 
 & HPD & Directed & 8,756 & 64,662 & 0 & 0.08 / 99.92 & 129,324 & 4,985,396 & 0.03 / 99.97 \\
 & YST & Directed & 2,284 & 13,292 & 0 & 0.25 / 99.75 & 26,584 & 562,721 & 0.08 / 99.92 \\
 & CEG & Directed & 297 & 4,296 & 0 & 4.89 / 95.11 & 8,592 & 305,630 & 0.41 / 99.59 \\
\midrule

% --- Social Networks ---
\multirow{3}{*}{\textbf{Social}} 
 & ADV & Directed & 5,155 & 78,570 & 0 & 0.30 / 99.70 & 157,140 & 14,420,896 & 0.06 / 99.94 \\
 & EML & Directed & 1,133 & 10,902 & 0 & 0.85 / 99.15 & 21,804 & 603,029 & 0.13 / 99.87 \\
 & BUP & Directed & 105 & 882 & 0 & 8.08 / 91.92 & 1,764 & 32,180 & 1.03 / 98.97 \\
\midrule

% --- Infrastructure / Transport ---
\multirow{2}{*}{\textbf{Infrastructure}} 
 & UAL & Directed & 332 & 4,252 & 0 & 3.87 / 96.13 & 8,504 & 417,821 & 0.58 / 99.42 \\
 & UPG & Directed & 4,941 & 13,188 & 0 & 0.05 / 99.95 & 26,376 & 168,349 & 0.02 / 99.98 \\

\bottomrule
\end{tabular}%
}
\end{table*}


\subsection{Statistical Significance Testing}
\label{subsec:statistical_tests}

To ensure robust evaluation of LineML's performance improvements, we employ two complementary statistical tests: the Wilcoxon signed-rank test for paired comparisons and Cliff's Delta for quantifying effect sizes. 

\subsubsection{Wilcoxon Signed-Rank Test}
The Wilcoxon signed-rank test \cite{wilcoxon1945individual} is a non-parametric statistical test used to compare paired samples when their differences cannot be assumed to follow a normal distribution. 
Unlike the Student's t-test, which assumes normality and equal variances, the Wilcoxon test makes no distributional assumptions, making it particularly suitable for performance metrics in machine learning where results often exhibit non-normal distributions.

In our context, for each dataset we treat LineML's performance scores (across multiple runs) as one sample and each baseline method's scores as another sample. 
The test works by:
\begin{enumerate}
    \item Calculating the differences between paired observations (LineML score minus baseline score for each experimental run)
    \item Ranking the absolute values of these differences, ignoring signs
    \item Summing the ranks for positive differences and negative differences separately
    \item Comparing the smaller sum to critical values to determine statistical significance
\end{enumerate}

A significant result (typically $p < 0.05$) indicates that LineML's performance is systematically different from the baseline method. 
Crucially, the test's non-parametric nature ensures validity even with small sample sizes (we use 5 runs per method) or when performance scores contain outliers.

\subsubsection{Cliff's Delta Effect Size}
While statistical significance tells us whether an effect exists, effect size quantifies its magnitude. 
Cliff's Delta $\delta$ \cite{cliff1993dominance} is a robust, non-parametric measure of effect size.
The test estimates the probability that a randomly selected score from one group exceeds a randomly selected score from another, minus the reverse probability. 
Formally, for two groups $X$ and $Y$:

\[
\delta = P(X > Y) - P(X < Y)
\]

where $\delta$ ranges from -1 to +1. Values are interpreted as:
\begin{itemize}[noitemsep,topsep=2pt]
    \item $|\delta| < 0.147$: Negligible effect
    \item $0.147 \leq |\delta| < 0.33$: Small effect
    \item $0.33 \leq |\delta| < 0.474$: Medium effect
    \item $|\delta| \geq 0.474$: Large effect
\end{itemize}

Cliff's Delta offers several advantages for our analysis. 
First, it is robust to outliers and non-normal distributions, complementing the Wilcoxon test's assumptions. 
Second, it provides intuitive interpretation: a $\delta$ of +0.6 means that in 80\% of random pairings, LineML outperforms the baseline method ($P(\text{LineML} > \text{Baseline}) = 0.8$). 
Finally, it quantifies practical significance by distinguishing between statistically significant but trivial differences versus meaningful improvements that matter in real applications.



\subsection{Functional Performance: Predictive Accuracy Assessment}


\begin{table}[ht!]
\centering
\caption{Link prediction performance (AUC \%) on citation networks. Results are averaged over \textbf{5 independent runs}. The top three performers are highlighted in \textbf{bold}, \textcolor{red}{red}, and \textcolor{blue}{blue}, respectively.}
\label{tab:citation_results_stacked}
\renewcommand{\arraystretch}{1.0}
\setlength{\tabcolsep}{3.5pt}
\small

\begin{tabular}{@{}lcccccr@{}}
\toprule

\multirow{2}{*}{Model} & \multicolumn{5}{c}{Dataset (AUC \%)} & \multirow{2}{*}{Avg. Rank} \\
\cmidrule(lr){2-6}
& NSC & ZWL & LDG & GRQ & Physics & \\
\midrule
LineML  & 99.51$\pm$0.27 & \textbf{97.84$\pm$3.06} & \textcolor{red}{96.86$\pm$1.72} & 96.35$\pm$0.76 & \textbf{97.70$\pm$0.04} & 2.4 \\
LGLP    & \textbf{99.82$\pm$0.01} & \textcolor{red}{97.76$\pm$0.01} & \textcolor{blue}{96.70$\pm$0.07} & \textcolor{blue}{97.68$\pm$0.10} & \textcolor{red}{97.25$\pm$0.78} & 2.7 \\
EG-DNMF    & 99.49$\pm$1.03 & 97.49$\pm$0.75 & 96.64$\pm$0.75 & \textcolor{red}{97.82$\pm$0.68} & \textcolor{blue}{97.15$\pm$0.53} & 4.4 \\
BSSLP   & \textcolor{blue}{99.63$\pm$0.69} & 97.45$\pm$0.57 & 96.52$\pm$0.82 & 96.36$\pm$0.13 & -- & 5.67 \\
DGI     & 99.19$\pm$0.32 & 96.41$\pm$1.02 & \textbf{97.49$\pm$0.39} & \textbf{98.30$\pm$0.37} & 96.15$\pm$0.42 & 5.8 \\
mLink   & \textcolor{red}{99.65$\pm$0.01} & \textcolor{blue}{97.67$\pm$0.01} & 96.62$\pm$0.11 & 97.56$\pm$0.10 & 95.78$\pm$0.42 & 6.0 \\
NRI     & 99.42$\pm$0.01 & 97.44$\pm$0.01 & 96.49$\pm$0.08 & 97.21$\pm$0.11 & 95.82$\pm$0.42 & 6.6 \\
SEAL    & 99.55$\pm$0.01 & 97.46$\pm$0.02 & 96.44$\pm$0.13 & 97.10$\pm$0.12 & 94.49$\pm$1.17 & 6.6 \\
GANE    & 99.36$\pm$0.95 & 83.97$\pm$1.18 & 96.39$\pm$0.51 & 97.22$\pm$2.15 & 95.55$\pm$0.75 & 8.45 \\
GAE     & 98.83$\pm$0.33 & 95.46$\pm$0.30 & 93.84$\pm$0.21 & 91.15$\pm$0.45 & 96.55$\pm$0.85 & 8.9 \\
Katz    & 98.00$\pm$0.31 & 96.42$\pm$0.12 & 92.96$\pm$0.19 & 89.81$\pm$0.59 & 95.56$\pm$0.15 & 11.65 \\
PR      & 98.05$\pm$0.29 & 97.20$\pm$0.12 & 94.46$\pm$0.19 & 89.98$\pm$0.57 & 73.36$\pm$0.18 & 11.9 \\
RB      & 96.95$\pm$0.04 & 95.18$\pm$0.03 & 90.00$\pm$0.10 & 90.88$\pm$0.01 & 94.18$\pm$0.03 & 12.3 \\
N2V     & 96.23$\pm$0.95 & 94.38$\pm$0.51 & 91.88$\pm$0.56 & 91.33$\pm$0.53 & 55.21$\pm$0.18 & 12.35 \\
SR      & 97.19$\pm$0.48 & 95.97$\pm$0.16 & 90.95$\pm$0.14 & 89.81$\pm$0.58 & 94.77$\pm$0.12 & 13.35 \\
\midrule


\multirow{2}{*}{Model} & \multicolumn{5}{c}{Dataset (AUC \%)} & \multirow{2}{*}{Avg. Rank} \\
\cmidrule(lr){2-6}
& SMG & KHN & PubMed & Cora & CiteSeer & \\
\midrule
LineML  & \textbf{95.31$\pm$0.40} & \textcolor{red}{95.02$\pm$0.11} & \textbf{98.85$\pm$0.08} & \textbf{95.14$\pm$0.58} & \textbf{96.85$\pm$0.08} & 2.4 \\
LGLP    & \textcolor{blue}{92.53$\pm$0.29} & 93.30$\pm$0.09 & \textcolor{red}{98.45$\pm$0.03} & \textcolor{blue}{94.45$\pm$0.03} & \textcolor{blue}{96.45$\pm$0.03} & 2.7 \\
EG-DNMF    & 91.21$\pm$0.72 & \textbf{95.15$\pm$0.84} & 97.41$\pm$0.87 & 94.41$\pm$0.87 & 96.41$\pm$0.65 & 4.4 \\
BSSLP   & 92.12$\pm$0.83 & \textcolor{blue}{94.43$\pm$0.66} & 95.28$\pm$0.36 & 93.28$\pm$0.36 & 96.05$\pm$0.36 & 5.67 \\
DGI     & \textcolor{red}{95.14$\pm$0.80} & 93.33$\pm$0.74 & 97.30$\pm$0.42 & 90.30$\pm$0.42 & 95.30$\pm$0.42 & 5.8 \\
mLink   & 92.05$\pm$0.32 & 92.89$\pm$0.11 & 96.27$\pm$0.23 & 92.27$\pm$0.23 & 92.27$\pm$0.23 & 6.0 \\
NRI     & 91.18$\pm$0.35 & 92.63$\pm$0.11 & 97.80$\pm$0.01 & 94.44$\pm$0.01 & 93.52$\pm$0.42 & 6.6 \\
SEAL    & 91.53$\pm$0.46 & 92.69$\pm$0.14 & 97.77$\pm$0.32 & 90.89$\pm$0.67 & \textcolor{red}{96.79$\pm$0.12} & 6.6 \\
GANE    & 91.56$\pm$0.97 & 80.18$\pm$0.33 & \textcolor{blue}{97.97$\pm$0.25} & 90.97$\pm$0.25 & 94.45$\pm$0.33 & 8.45 \\
GAE     & 85.88$\pm$0.90 & 84.37$\pm$0.39 & 97.45$\pm$0.07 & \textcolor{red}{95.08$\pm$0.27} & 93.45$\pm$0.47 & 8.9 \\
Katz    & 86.09$\pm$1.06 & 84.60$\pm$0.79 & 74.86$\pm$0.52 & 81.17$\pm$0.52 & 75.86$\pm$0.58 & 11.65 \\
PR      & 89.13$\pm$0.90 & 88.43$\pm$0.80 & 63.20$\pm$0.01 & 70.20$\pm$0.01 & 64.20$\pm$0.52 & 11.9 \\
RB      & 86.37$\pm$0.50 & 80.24$\pm$0.17 & 93.45$\pm$0.03 & 90.42$\pm$0.12 & 92.41$\pm$0.85 & 12.3 \\
N2V     & 78.30$\pm$1.22 & 82.21$\pm$1.19 & 93.14$\pm$0.80 & 90.97$\pm$0.64 & 93.14$\pm$0.25 & 12.35 \\
SR      & 78.39$\pm$1.14 & 79.55$\pm$0.90 & 63.97$\pm$0.56 & 73.97$\pm$0.45 & 68.97$\pm$0.45 & 13.35 \\

\bottomrule
\end{tabular}
\end{table}
\newpage
\begin{table}[ht!]
\centering
\caption{ink prediction performance (AUC \%) on citation networks. Results are averaged over \textbf{5 independent runs}. The top three performers are highlighted in \textbf{bold}, \textcolor{red}{red}, and \textcolor{blue}{blue}, respectively.}
\label{tab:social_bio_results_split}
\renewcommand{\arraystretch}{1.0}
\setlength{\tabcolsep}{3.5pt}
\small
\begin{tabular}{@{}lcccccr@{}}
\toprule

\multirow{2}{*}{Model} & \multicolumn{4}{c}{Dataset (AUC \%)} & \multirow{2}{*}{Avg. Rank} \\
\cmidrule(lr){2-5}
& UAL & ADV & EML & BUP & \\
\midrule
LineML & \textbf{98.13$\pm$0.14} & \textbf{95.51$\pm$0.06} & \textbf{94.55$\pm$0.09} & \textbf{98.78$\pm$0.01} & 1 \\
LGLP & \textcolor{blue}{97.44$\pm$0.32} & \textcolor{red}{95.40$\pm$0.10} & 92.03$\pm$0.28 & \textcolor{red}{95.24$\pm$0.53} & 3.94 \\
EG-DNMF & \textcolor{red}{97.83$\pm$0.41} & 94.89$\pm$0.19 & 92.14$\pm$0.41 & 85.59$\pm$0.79 & 4.88 \\
mLink & 96.43$\pm$0.33 & 95.21$\pm$0.10 & 92.03$\pm$0.31 & 93.54$\pm$0.63 & 5.44 \\
GANE & 95.18$\pm$0.57 & \textcolor{blue}{95.31$\pm$0.60} & \textcolor{blue}{93.74$\pm$5.46} & \textcolor{blue}{94.97$\pm$2.65} & 5.75 \\
NRI & 96.44$\pm$0.38 & 94.53$\pm$0.10 & 91.83$\pm$0.30 & 94.77$\pm$0.60 & 6.25 \\
BSSLP & 97.02$\pm$0.75 & 95.20$\pm$0.26 & 91.68$\pm$0.82 & 77.24$\pm$0.53 & 6.12 \\
DGI & 88.63$\pm$0.34 & 91.37$\pm$2.24 & \textcolor{red}{94.25$\pm$0.54} & 93.74$\pm$1.43 & 7 \\
SEAL & 95.21$\pm$0.77 & 95.07$\pm$0.13 & 92.01$\pm$0.38 & 93.32$\pm$0.84 & 7.62 \\
PR & 93.74$\pm$1.01 & 92.78$\pm$0.18 & 89.46$\pm$0.63 & 90.13$\pm$2.45 & 9.5\\
Katz & 92.01$\pm$0.88 & 92.13$\pm$0.21 & 88.45$\pm$0.68 & 87.10$\pm$2.73 & 11.25 \\
RB & 90.25$\pm$0.54 & 90.50$\pm$0.04 & 87.12$\pm$0.29 & 87.20$\pm$1.01 & 11.38 \\
GAE & 91.80$\pm$0.86 & 90.55$\pm$0.23 & 86.78$\pm$1.07 & 90.16$\pm$1.65 & 12.06 \\
SR & 79.21$\pm$1.50 & 86.18$\pm$0.22 & 86.90$\pm$0.71 & 85.47$\pm$2.75 & 13.88 \\
N2V & 85.40$\pm$0.96 & 77.70$\pm$0.83 & 83.06$\pm$1.42 & 80.25$\pm$5.55 & 13.97 \\

\end{tabular}

\begin{tabular}{@{}lcccccr@{}}
\toprule
\multirow{2}{*}{Model} & \multicolumn{4}{c}{Dataset (AUC \%)} & \multirow{2}{*}{Avg. Rank} \\
\cmidrule(lr){2-5}
& HPD & CEG & YST & UPG & \\
\midrule
LineML & \textbf{96.48$\pm$1.66} & \textbf{96.01$\pm$0.09} & \textbf{94.39$\pm$0.13} & \textbf{94.53$\pm$0.36} & 1 \\
LGLP & 92.58$\pm$0.08 & 90.16$\pm$0.76 & \textcolor{red}{91.97$\pm$0.12} & 82.17$\pm$0.57 & 3.94 \\
EG-DNMF & \textcolor{blue}{94.49$\pm$0.37} & \textcolor{red}{92.89$\pm$1.74} & 91.26$\pm$0.87 & 89.53$\pm$0.31 & 4.88 \\
mLink & 92.64$\pm$0.08 & 89.08$\pm$0.86 & 91.40$\pm$0.13 & 83.14$\pm$0.61 & 5.44 \\
GANE & 92.46$\pm$0.82 & 88.79$\pm$2.87 & 78.56$\pm$0.48 & \textcolor{red}{92.31$\pm$0.57} & 5.75 \\
NRI & 91.63$\pm$0.08 & 90.13$\pm$0.82 & \textcolor{blue}{91.63$\pm$0.25} & 82.15$\pm$0.50 & 6.12 \\
BSSLP & 93.76$\pm$0.52 & \textcolor{blue}{91.53$\pm$0.77} & 90.54$\pm$1.18 & \textcolor{blue}{90.92$\pm$0.05} & 6.25 \\
DGI & \textcolor{red}{96.01$\pm$0.82} & 85.77$\pm$1.70 & 85.03$\pm$0.55 & 83.43$\pm$0.16 & 7 \\
SEAL & 91.71$\pm$0.25 & 87.44$\pm$1.21 & 82.07$\pm$0.96 & 81.37$\pm$0.93 & 7.62 \\
PR & 87.19$\pm$0.34 & 89.14$\pm$1.35 & 81.40$\pm$0.75 & 59.88$\pm$1.51 & 9.5 \\
Katz & 85.47$\pm$0.35 & 84.84$\pm$2.05 & 80.56$\pm$0.78 & 59.59$\pm$1.51 & 11.25 \\
RB & 82.89$\pm$0.09 & 85.94$\pm$1.89 & 80.39$\pm$0.08 & 75.93$\pm$0.08 & 11.38 \\
GAE & 85.21$\pm$0.45 & 83.73$\pm$0.75 & 77.07$\pm$0.36 & 69.84$\pm$0.96 & 12.06 \\
SR & 81.73$\pm$0.37 & 75.65$\pm$2.24 & 73.93$\pm$0.95 & 70.18$\pm$0.75 & 13.88\\
N2V & 79.61$\pm$1.14 & 80.08$\pm$1.52 & 77.07$\pm$0.36 & 70.37$\pm$1.15 & 13.94 \\
\bottomrule
\end{tabular}
\end{table}

\newpage


\subsubsection{Evaluation Methodology and Baseline Configurations}
\label{subsub:Basline}
To evaluate LineML framework, we perform a comparative analysis against a diverse suite of 14 benchmarking techniques across 18 datasets presented in table \ref{tab:dataset}. 
For clarity, we first introduce the rationale and a brief description of each baseline method.
Our evaluation encompasses the following methods, grouped by their underlying paradigm:
\begin{itemize}
    \item \textbf{Classical Heuristics:} These topology-driven methods provide a fundamental performance baseline. 
    We include the \textit{Katz} index \cite{24katz1953new}, which sums over all paths between nodes with an attenuation factor; 
    \textit{PageRank (PR)} \cite{25brin2012reprint}, adapted for link prediction via rooted random walks; and 
    \textit{SimRank (SR)} \cite{26jeh2002simrank}, which measures similarity based on the context of interconnected nodes. 
    We also evaluate two recent, high-performing heuristics: 
    the \textit{Resource Broadcast (RB)} index \cite{liu2024rb}, which refines resource transmission logic, and 
    \textit{BSSLP} \cite{liu2025bsslp}, a hybrid method designed to resolve the zero-similarity problem in sparse networks.
  
    \item \textbf{Embedding-Based Approaches:} This category includes methods that learn low-dimensional node representations. 
    We compare against \textit{Node2vec (N2V)} \cite{40grover2016node2vec}, which learns embeddings through biased random walks, and the \textit{Graph Auto-Encoder (GAE)} \cite{31kipf2016variational}, which reconstructs the adjacency matrix via a variational autoencoder architecture.
    
    \item \textbf{GNN-Based Approaches:} We benchmark against several influential neural models. 
    This includes \textit{SEAL} \cite{14zhang2018link}, which formulates link prediction as a subgraph classification task;
    \textit{Multi-Scale Link Prediction (mLink)} \cite{33cai2020multi}, an extension of SEAL that captures multi-scale subgraph features; 
    \textit{Deep Graph Infomax (DGI)} \cite{velickovic2018deep}, 
    a self-supervised method based on mutual information maximization; 
    \textit{Generative Adversarial Network Embedding (GANE)} \cite{hong2020gane}, which employs adversarial training for robust embeddings; and \textit{Neural Relational Inference (NRI)} \cite{41kipf2018neural}, a model that infers latent interaction graphs.
    
    \item \textbf{Matrix Factorization Methods:} We include the deep non-negative matrix factorization approach \textit{EG-DNMF} \cite{yao2024deep}, notable for its integrated edge generator to handle network sparsity.
    
    \item \textbf{Line Graph Transformation:} As a direct precursor to our work, we evaluate \textit{Line Graph-based Link Prediction (LGLP)} \cite{cai2021line}, which first established the link-to-node classification reformulation via line graphs.
\end{itemize}

\paragraph{Evaluation Protocol and Parameter Settings}
To ensure a fair and reproducible comparison, all baseline methods were evaluated following the protocols and parameter ranges recommended in their original publications. 
For classical heuristic approaches, we adopt commonly used configurations: the damping factor of Katz is set to \(0.001\) \cite{cai2021line,33cai2020multi}, 
PageRank uses a damping factor of \(0.85\), and 
SimRank employs a decay coefficient of \(0.8\). 
For learning-based baselines, including Node2vec, GAE, SEAL, mLink, DGI, GANE, NRI, EG-DNMF, and LGLP, we perform a lightweight hyperparameter search over the ranges reported in the corresponding papers, selecting the best configuration based on validation performance. 
We use 80\% of edges for training, 10\% for validation, and 10\% for testing. 

The hyperparameters of the proposed LineML framework are similarly determined using validation data. The adaptive margin \(\alpha\) in the triplet loss is fixed to \(1.0\). Optimization is performed using Adam, with learning rates of \(1\times10^{-4}\) for the GraphSAGE encoder and \(1\times10^{-3}\) for the link classifier. This consistent validation-driven tuning strategy ensures a balanced comparison between LineML and all competing approaches.
Negative edges are  degree-based sampled from non-existent edges, maintaining a 1:1 positive-to-negative ratio. Experiments are repeated over 5 random seeds to ensure statistical reliability.


\paragraph{Evaluation Metric and Statistical Validation}
We employ the Area Under the Receiver Operating Characteristic Curve (AUC) as the primary evaluation metric. 
AUC measures the ability of a model to rank positive links higher than negative links, providing a threshold-independent assessment of discriminative performance.

To assess statistical significance, we rely on the Wilcoxon Signed-Rank test in conjunction with Cliff’s delta effect size. 
These complementary non-parametric and robust measures enable reliable paired comparisons across datasets and provide both statistical and practical significance. 









\subsubsection{Cross-Method Performance Analysis Across Network Categories}
\label{subsec:comparative_analysis}



Tables~\ref{tab:citation_results_stacked} and~\ref{tab:social_bio_results_split} present a comprehensive performance evaluation of the LineML method against the 14 baseline models across the 18 datasets. 
Results are reported using the AUC metric, denoted as the mean $\pm$ standard deviation calculated over five independent runs. 
For clarity, the datasets are partitioned into two primary categories: citation networks (Table~\ref{tab:citation_results_stacked}) and social/biological networks (Table~\ref{tab:social_bio_results_split}).




\paragraph{Performance Analysis in Citation Networks} 
In citation networks (Table~\ref{tab:citation_results_stacked}), LineML remains highly competitive, achieving the best or second-best performance in 8 of 10 datasets. 
It attains the highest average rank (2.4), slightly ahead of LGLP (2.7). 
The margin of superiority varies with network density and regularity. 
In dense, well-structured citation graphs like NSC and GRQ, performance differences become negligible; LineML ($99.51\%$), LGLP ($99.82\%$), and EG-DNMF ($97.76\%$) all achieve excellent results. 
This suggests that in uniform networks with clear citation patterns, even simpler graph representations can capture most relevant information.
Conversely, in more challenging citation graphs like Physics and ZWL, LineML's advantages resurface. 
It achieves $97.70\%$ on Physics, outperforming LGLP ($97.25\%$) and EG-DNMF ($97.15\%$). 
The Physics graph is notably large and complex, indicating that LineML's adaptive metric learning provides greater robustness to scale and noise. 
This pattern reveals an important insight: {the value of LineML's architecture correlates with network complexity}. 
In simpler, denser graphs, its benefits are marginal, but in larger, noisier, or more heterogeneous networks, its advantages become pronounced.
Wilcoxon signed-rank test shows that LineML maintains {strong superiority}, with perfect wins against classical methods (N2V, Katz, PR, SR, RB) and significant advantages over embedding-based approaches (GAE, GANE: $p < 0.01$). Against advanced GNN competitors (DGI, NRI, SEAL, mLink), LineML wins 80\% of comparisons ($p < 0.05$), though the closest competition comes from other line graph methods: LGLP and EG-DNMF remain competitive with 30\% and 20\% win rates respectively against LineML.

\paragraph{Performance Analysis in Social/Biological Networks}
LineML demonstrates particularly strong performance in social and biological networks, achieving top-rank AUC scores in all eight datasets of this category (Table~\ref{tab:social_bio_results_split}). 
The improvement over alternatives is often substantial and statistically significant. 
For instance, in the BUP political network, LineML achieves $98.78 $ AUC compared to LGLP's $95.24 $. 
Similarly, in the UPG protein interaction network, LineML scores $94.53$ versus LGLP's $82.17$. 
These networks typically exhibit higher structural heterogeneity and complex relational patterns beyond simple homophily. 
LineML's adaptive triplet loss and line graph transformation prove especially effective here.
They dynamically adjust to varying similarity thresholds and capture nuanced edge relationships that other methods miss.

\paragraph{Methodological Category Analysis}
Analyzing performance by methodological category reveals consistent trends:
Classical and Heuristic Methods (Katz, PR, AA, PA): These methods consistently underperform, particularly in social/biological networks. 
In the ADV social network, for example, Katz achieves $92.13\%$ compared to LineML's $95.51\%$. 
Their reliance on fixed structural formulas limits adaptability to diverse relationship types.
Embedding Methods (GAE, N2V, DGI): Performance varies significantly. While DGI occasionally competes (tying for second in HPD with $96.01\%$), methods like N2V struggle, especially in sparse networks ($80.08\%$ in CEG vs. LineML's $96.01\%$). 
Embedding methods often fail to explicitly model relational dependencies between edges, a core strength of line graph approaches.
{GNN-Based Competitors (SEAL, mLink, NRI)}: These represent the strongest competitors, particularly in citation networks. 
In the GRQ dataset, mLink achieves $97.56\%$, close to LineML's $96.35\%$. 
However, in more heterogeneous networks like YST, the gap widens: LineML ($94.39\%$) outperforms mLink ($91.40\%$) and SEAL ($82.07\%$). 
This suggests that while GNN-based methods effectively aggregate local neighborhood information, LineML's explicit edge-to-edge modeling captures relational patterns more effectively.
As the most similar approaches, these provide the closest competition. 
LGLP performs strongly in citation networks, even topping LineML in NSC and GRQ. 
However, LineML's integration of adaptive metric learning consistently provides an edge in more complex scenarios, as evidenced by its superior average rank. 
The performance gap is most significant in networks with high relational heterogeneity, such as BUP and UPG.



\paragraph{Statistical Significance Validation}
Figure~\ref{fig:cliffs_delta_barplot} quantifies LineML's performance advantage through Cliff's Delta effect sizes, revealing distinct patterns across network types. 
In social and biological networks, LineML achieves near-perfect dominance ($\delta \geq 0.97$) against all baselines, demonstrating superiority in heterogeneous relational environments.
In citation networks, the competitive landscape sharpens. 
While LineML maintains large effects against classical methods ($\delta \geq 0.96$), advanced approaches provide stronger competition: LGLP ($\delta = 0.168$) and EG-DNMF ($\delta = 0.32$) show small-to-medium effects, reflecting their capacity to capture structural patterns in well-organized academic graphs.
Globally, LineML achieves large effects ($\delta \geq 0.73$) against most methods, with LGLP ($\delta = 0.52$) and EG-DNMF ($\delta = 0.57$) representing the closest competitors. 


\begin{figure}[!ht]
    \centering
    \includegraphics[width=0.95\textwidth]{figs/cliffs_delta_summary_refined.pdf}
    \caption{Cliff's Delta effect sizes for LineML versus baseline methods. Values approaching +1.0 indicate near-complete dominance.}
    \label{fig:cliffs_delta_barplot}
\end{figure}




\begin{figure*}[ht!]
    \centering
    \begin{minipage}[b]{0.30\textwidth}
        \includegraphics[width=\textwidth]{figs/2-BUP_ROC_errorbar.pdf}
    \end{minipage}
    \begin{minipage}[b]{0.30\textwidth}
        \includegraphics[width=\textwidth]{figs/2-ADV_ROC_errorbar.pdf}
    \end{minipage}
    \begin{minipage}[b]{0.30\textwidth}
        \includegraphics[width=\textwidth]{figs/2-CEG_ROC_errorbar.pdf}
    \end{minipage}

    \vspace{0.2cm}

    \begin{minipage}[b]{0.30\textwidth}
        \includegraphics[width=\textwidth]{figs/2-EML_ROC_errorbar.pdf}
    \end{minipage}
    \begin{minipage}[b]{0.30\textwidth}
        \includegraphics[width=\textwidth]{figs/2-GRQ_ROC_errorbar.pdf}
    \end{minipage}
    \begin{minipage}[b]{0.30\textwidth}
        \includegraphics[width=\textwidth]{figs/2-HPD_ROC_errorbar.pdf}
    \end{minipage}

    \vspace{0.2cm}

    \begin{minipage}[b]{0.30\textwidth}
        \includegraphics[width=\textwidth]{figs/2-KHN_ROC_errorbar.pdf}
    \end{minipage}
    \begin{minipage}[b]{0.30\textwidth}
        \includegraphics[width=\textwidth]{figs/2-LDG_ROC_errorbar.pdf}
    \end{minipage}
    \begin{minipage}[b]{0.30\textwidth}
        \includegraphics[width=\textwidth]{figs/2-NSC_ROC_errorbar.pdf}
    \end{minipage}

    \vspace{0.2cm}

    \begin{minipage}[b]{0.30\textwidth}
        \includegraphics[width=\textwidth]{figs/2-SMG_ROC_errorbar.pdf}
    \end{minipage}
    \begin{minipage}[b]{0.30\textwidth}
        \includegraphics[width=\textwidth]{figs/2-UAL_ROC_errorbar.pdf}
    \end{minipage}
    \begin{minipage}[b]{0.30\textwidth}
        \includegraphics[width=\textwidth]{figs/2-UPG_ROC_errorbar.pdf}
    \end{minipage}

    \vspace{0.2cm}

    \begin{minipage}[b]{0.30\textwidth}
        \includegraphics[width=\textwidth]{figs/2-YST_ROC_errorbar.pdf}
    \end{minipage}
    \begin{minipage}[b]{0.30\textwidth}
        \includegraphics[width=\textwidth]{figs/2-ZWL_ROC_errorbar.pdf}
    \end{minipage}
    
    \caption{Comparison of AUC scores across multiple datasets for the proposed LineML method and baseline methods (Katz, PR, SR, SEAL, and LGLP) using varying proportions of training links. }
\label{fig:training_size_effect}
\end{figure*}







\subsubsection{Training Size Sensitivity}

In Figure~\ref{fig:training_size_effect}, we analyze AUC performance across multiple datasets for the proposed LineML framework and representative baseline methods (Katz, PR, SR, SEAL, and LGLP) under varying proportions of training links. These baselines were selected to cover complementary families of link prediction approaches: Katz represents classical path-based similarity, PR and SR capture random-walk and neighborhood-driven heuristics, SEAL serves as a strong subgraph-based GNN model, and LGLP reflects recent line graph-based learning strategies. Together, they provide a balanced comparison between traditional heuristics, modern GNN architectures, and line graph paradigms. For each dataset, six training ratios are evaluated (30\%, 40\%, 50\%, 60\%, 70\%, and 80\%). 
LineML is depicted by a solid black curve with shaded standard deviation, while baseline methods are illustrated using dashed lines, enabling a clear visualization of relative trends and variability.

Across datasets such as CEG, BUP, HPD, EML, KHN, LDG, SMG, and UPG, LineML consistently outperforms all baselines by a substantial margin over the full range of training sizes. 
This behavior indicates strong robustness to reduced supervision and highlights LineML’s ability to generalize effectively in structurally diverse graphs. 
The observed performance gap can be attributed to LineML’s joint exploitation of line graph transformation and metric learning, which enables the model to capture higher-order structural dependencies that are difficult to represent using classical similarity measures or GNNs relying primarily on local neighborhood aggregation.

For the ADV and YST datasets, LineML maintains a consistent advantage, although with a narrower margin. 
This suggests that these graphs may exhibit simpler structural patterns, allowing baseline methods to recover a considerable portion of link information without requiring advanced representation learning. 
The reduced gap indicates that neighborhood-based or path-based signals are already informative in these settings. 
Nevertheless, LineML’s continued superiority demonstrates its adaptability and confirms that its design remains beneficial even when graph structures are relatively straightforward.

In GRQ, NSC, UAL, and ZWL, LineML shows slightly lower performance than certain baselines at smaller training ratios (30\% and 40\%), while remaining competitive overall. 
This behavior suggests that LineML may require a minimum amount of supervision to fully exploit its representation capacity, particularly in graphs characterized by sparsity or irregular connectivity. Under limited data, simpler baselines based on direct aggregation or heuristic similarity can achieve comparable results, as complex feature learning has not yet converged. As training size increases, LineML consistently closes this gap and typically surpasses competing methods, underscoring its effectiveness once sufficient training evidence becomes available.








\subsection{Non-Functional Performance: Computational Efficiency Analysis}
This subsection presents a comprehensive evaluation of the computational efficiency and scalability of our proposed LineML framework. 
While predictive accuracy (functional performance) is essential for link prediction models, practical deployment in real-world applications requires careful consideration.
We systematically assess three key dimensions.
First, the effectiveness of line graph pruning strategies in reducing computational overhead while preserving accuracy. 
Second, the performance benefits of distributed training implementations. 
Finally, comparative computational efficiency against state-of-the-art baseline methods. 
Through experiments across small graphs (e.g., Cora, CiteSeer) to large co-authorship graphs (e.g., Physics), we quantify the trade-offs between accuracy and efficiency.
To isolate pruning effects, all experiments use fixed optimal hyperparameters (Section~\ref{sec:SysAblation}): $(\alpha, \beta) = (1.0, 1.0)$, mean GraphSAGE aggregator, 1:1 sampling ratio, batch size 256, classifier learning rate $10^{-6}$, GNN learning rate $0.001$, three GNN layers, three classifier layers, and four attention heads.


\subsubsection{Pruning Strategies: Performance-Scalability Trade-offs}
\label{subsec:pruning_results}



We compare the three pruning approaches: Degree-Based Hierarchical Pruning, Spectral Sparsification via Effective Resistance, and kNN Spectral Pruning. 
For degree pruning, we evaluate both \textit{high-degree} and \textit{low-degree} removal modes with thresholds chosen to obtain multiple pruning ratios (ranging from approximately 20\% to 80\% edge removal). 
For spectral sparsification, we vary the preservation ratio $\lambda \in \{0.1, 0.3, 0.5, 0.7, 0.9\}$, directly controlling the fraction of retained edges. 
For kNN pruning, we test $k \in \{2,5,10\}$, where each node keeps its $k$ nearest neighbors in feature space. 

Table~\ref{tab:pruning_results} reports AUC together with training time per epoch. 
The column ``Pruned (\%)'' indicates the fraction of removed edges, while ``Imp'' denotes the relative change in AUC compared to the unpruned baseline.


Across all datasets, aggressive pruning generally leads to performance degradation, especially when more than 80--90\% of edges are removed. 
This effect is most visible for spectral pruning at $\lambda=0.1$ and for kNN with $k=2$, where excessive sparsification disrupts local neighborhood structure and weakens message passing. 
In contrast, moderate pruning ratios often preserve, and in some cases slightly improve, predictive performance. 
This behavior suggests that pruning can act as a form of regularization by removing redundant or noisy connections.

Among all methods, kNN pruning with $k=10$ consistently achieves the best trade-off. 
Despite removing approximately 76\% of line graph edges, it maintains or improves AUC on all datasets, including large-scale Physics, where it even yields a $+0.38\%$ gain. 
This confirms that preserving feature-consistent neighborhoods is sufficient to retain discriminative structure for link prediction, while discarding a large number of less informative connections.

Degree-based pruning exhibits mixed behavior. 
While high-degree pruning can maintain performance at moderate ratios, low-degree pruning often leads to noticeable degradation, particularly on CiteSeer and Physics. 
This indicates that peripheral nodes still contribute useful information and that removing them too aggressively harms global representation quality. 
Spectral pruning provides strong theoretical guarantees, but in practice requires relatively large preservation ratios ($\lambda \geq 0.7$) to avoid performance loss, limiting its compression benefits.


As expected from Lemma~\ref{lemma:trainingTimeReduction}, pruning substantially reduces training time per epoch, with speedups roughly proportional to the pruning ratio. 
On small graphs such as Cora and CiteSeer, gains are moderate (typically $1.2\times$--$2.4\times$), reflecting limited absolute workload. 
However, on larger networks the impact becomes pronounced. On Physics, kNN pruning with $k=10$ accelerates training by up to $4.6\times$, reducing per-epoch time from 1,793\,s to 389\,s, while preserving accuracy. Even more aggressive configurations can reach speedups above $10\times$, albeit at the cost of performance.

These results highlight that pruning is especially beneficial for large-scale graphs, where edge explosion in the line graph dominates both computation and communication. 
By reducing adjacency density, pruning improves cache locality and lowers inter-partition message passing in distributed environments, making it particularly suitable for HPC settings.



\begin{table}[ht]
\centering
\caption{Comparison of Pruning Methods. Bold values indicate the best results.}
\label{tab:pruning_results}

% Adjust font size and column spacing to fit width without resizebox
\footnotesize 
\setlength{\tabcolsep}{3.5pt} 
\renewcommand{\arraystretch}{1.1}

\begin{tabular}{llccccc}
\toprule

\multirow{2}{*}{\textbf{Methods}} & \multirow{2}{*}{\textbf{Param}} & \multirow{2}{*}{\textbf{Pruned (\%)}} & \multicolumn{2}{c}{\textbf{Cora}} & \multicolumn{2}{c}{\textbf{CiteSeer}} \\
\cmidrule(lr){4-5} \cmidrule(lr){6-7} 
 & & & \textbf{AUC (Imp)} & \textbf{Time} & \textbf{AUC (Imp)} & \textbf{Time (s)/epochs} \\
\midrule
\multicolumn{3}{l}{\textit{Baseline (Unpruned)}} & 95.14 & 10.79s & 96.85 & 13.29s \\
\midrule
\multirow{6}{*}{Degree Pruning} 
 &  \textbf{high $\tau =0.1$}& \textbf{64.1\%} & \textbf{95.19 (0.05\%)} & 8.8s (1.2$\times$) & \textbf{95.03 (-1.88\%)} & 6.4s (2.1$\times$) \\
 & low  $\tau =0.1$ & 35.9\% & 95.22 (0.08\%) & 9.0s (1.2$\times$) & 90.56 (-6.49\%) & 8.2s (1.6$\times$) \\
 & high $\tau =0.3$ & 31.2\% & 95.06 (-0.08\%) & 9.4s (1.1$\times$) & 94.04 (-2.90\%) & 7.9s (1.7$\times$) \\
 & low $\tau  =0.3$& 68.8\% & 93.04 (-2.21\%) & 7.4s (1.5$\times$) & 89.79 (-7.29\%) & 6.9s (1.9$\times$) \\
 & high $\tau  =0.5$ & 19.3\% & 95.84 (0.74\%) & 15.2s (0.7$\times$) & 94.33 (-2.60\%) & 9.7s (1.4$\times$) \\
 & low $\tau  =0.5$& 80.7\% & 90.85 (-4.51\%) & 6.2s (1.7$\times$) & 88.71 (-8.40\%) & 6.4s (2.1$\times$) \\
\midrule
\multirow{5}{*}{Spectral Pruning} 
 & $\lambda=$0.1 & 90.0\% & 70.91 (-25.47\%) & 7.4s (1.5$\times$) & 76.98 (-20.52\%) & 6.7s (2.0$\times$) \\
 & $\lambda=$0.3 & 70.0\% & 82.20 (-13.60\%) & 8.2s (1.3$\times$) & 89.70 (-7.38\%) & 8.5s (1.6$\times$) \\
 & $\lambda=$0.5 & 50.0\% & 90.58 (-4.79\%) & 8.6s (1.3$\times$) & 91.81 (-5.20\%) & 11.5s (1.2$\times$) \\
 & \textbf{$\lambda=$0.7} & \textbf{30.0\%} & \textbf{94.39 (-0.79\%)} & \textbf{9.9s (1.1$\times$)} & \textbf{92.40 (-4.59\%)} & \textbf{12.9s (1.0$\times$)} \\
 & $\lambda=$0.9 & 10.0\% & 95.06 (-0.08\%) & 10.0s (1.1$\times$) & 93.30 (-3.67\%) & 13.0s (1.0$\times$) \\
\midrule
\multirow{3}{*}{Knn Pruning} 
 & $k=$2 & 94.9\% & 85.62 (-10.01\%) & 6.5s (1.7$\times$) & 88.61 (-8.51\%) & 5.5s (2.4$\times$) \\
 & $k=$5 & 87.8\% & 93.18 (-2.06\%) & 6.2s (1.7$\times$) & 93.69 (-3.26\%) & 6.5s (2.0$\times$) \\
 & \textbf{$k=$10} & \textbf{75.9\%} & \textbf{95.61 (0.49\%)} & \textbf{8.8s (1.2$\times$)} & \textbf{96.01 (-0.87\%)} & \textbf{6.7s (2.0$\times$)} \\
\bottomrule


\multirow{2}{*}{\textbf{}} & \multirow{2}{*}{\textbf{}} & \multirow{2}{*}{\textbf{}} & \multicolumn{2}{c}{\textbf{PubMed}} & \multicolumn{2}{c}{\textbf{Physics}} \\
\cmidrule(lr){4-5} \cmidrule(lr){6-7} 


\multicolumn{3}{l}{\textit{Baseline (Unpruned)}} & 98.85 & 119.84s & 97.70 & 1,793.08s \\
\midrule
\multirow{6}{*}{Degree Pruning} 
 & \textbf{high $\tau =0.1$} & \textbf{64.1\%} & \textbf{98.93 (0.08\%)} & \textbf{106.3s (1.1$\times$)} & \textbf{97.60 (-0.10\%)} & 730.8s (2.5$\times$) \\
 & low $\tau =0.1$& 35.9\% & 98.37 (-0.49\%) & 111.2s (1.1$\times$) & 97.21 (-0.50\%) & 1082.4s (1.7$\times$) \\
 & high$\tau =0.3$ & 31.2\% & 98.92 (0.07\%) & 118.3s (1.0$\times$) & 97.41 (-0.30\%) & 1289.6s (1.4$\times$) \\
 & low $\tau =0.3$& 68.8\% & 92.56 (-6.36\%) & 79.2s (1.5$\times$) & 93.42 (-4.38\%) & 255.0s (7.0$\times$) \\
 & high $\tau =0.5$& 19.3\% & 99.09 (0.24\%) & 118.2s (1.0$\times$) & 97.42 (-0.29\%) & 1617.6s (1.1$\times$) \\
 & low $\tau =0.5$& 80.7\% & 93.43 (-5.48\%) & 71.5s (1.7$\times$) & 91.91 (-5.93\%) & 164.5s (10.9$\times$) \\
\midrule
\multirow{5}{*}{Spectral Pruning} 
 & $\lambda=$0.1 & 90.0\% & 97.51 (-1.36\%) & 76.4s (1.6$\times$) & 85.05 (-12.95\%) & 708.9s (2.5$\times$) \\
 & $\lambda=$0.3 & 70.0\% & 98.17 (-0.69\%) & 80.5s (1.5$\times$) & 92.36 (-5.47\%) & 1127.5s (1.6$\times$) \\
 & $\lambda=$0.5 & 50.0\% & 98.63 (-0.22\%) & 101.8s (1.2$\times$) & 95.43 (-2.32\%) & 1194.7s (1.5$\times$) \\
 & \textbf{$\lambda=$0.7} & \textbf{30.0\%} & \textbf{98.88 (0.03\%)} & \textbf{116.9s (1.0$\times$)} & \textbf{97.43 (-0.28\%)} & \textbf{1520.3s (1.2$\times$)} \\
 & $\lambda=$0.9 & 10.0\% & 98.90 (0.05\%) & 118.1s (1.0$\times$) & 97.54 (-0.16\%) & 1621.3s (1.1$\times$) \\
\midrule
\multirow{3}{*}{Knn Pruning} 
 & $k=$2 & 94.9\% & 94.69 (-4.21\%) & 76.7s (1.6$\times$) & 94.80 (-2.97\%) & 167.3s (10.7$\times$) \\
 & $k=$5 & 87.8\% & 96.51 (-2.37\%) & 77.7s (1.5$\times$) & 95.03 (-2.73\%) & 231.0s (7.8$\times$) \\
 & \textbf{$k=$10} & \textbf{75.9\%} & \textbf{98.94 (0.09\%)} & \textbf{88.3s (1.4$\times$)} & \textbf{98.07 (0.38\%)} & \textbf{389.4s (4.6$\times$)} \\
\bottomrule
\end{tabular}
\end{table}









\subsubsection{High-Performance Computing Integration and Scalability}
\label{subsec:distributed_hpc}


To evaluate our pruning framework in HPC environments, we implemented a distributed pipeline and tested eight configurations across four benchmark datasets. 
Each pruning method uses its previously identified optimal parameter: kNN with $k=10$, Degree with high pruning ($\tau=0.1$), and Spectral with $\lambda=0.7$. 
The tested configurations are:
\begin{itemize}
    \item \textbf{Sequential}: Baseline LineML; LineML with kNN pruning; LineML with Degree pruning; LineML with Spectral pruning
    \item \textbf{Parallel}: Parallel LineML; Parallel LineML with kNN pruning; Parallel LineML with Degree pruning; Parallel LineML with Spectral pruning
\end{itemize}

This design isolates the computational benefits of pruning, parallelization, and their combined effects on training efficiency and scalability.
All experiments are conducted in a HPC setting using PyTorch Distributed Data Parallel in a data-parallel configuration across two GPUs, with one training process per GPU and NCCL-based all-reduce for synchronized gradient updates. 
Mixed-precision training is enabled via automatic mixed precision with gradient scaling to improve computational efficiency. 
The full line graph is first constructed and pruned centrally on the CPU, then placed into shared memory and accessed concurrently by all worker processes without duplication. 
Rather than partitioning the graph topology, the global set of training node indices is evenly split across GPUs, while each worker performs local neighborhood sampling using PyTorch Geometric’s \texttt{NeighborLoader}, preserving the complete graph structure on each device. 
Evaluation is performed after synchronized training. 
This combination of multi-GPU parallelism, shared-memory large-scale graph processing, and CUDA-accelerated mixed-precision training enables efficient execution beyond single-device capabilities, justifying the use of an HPC environment.



Figure~\ref{fig:fused_pruning_comparison} presents the comprehensive speedup analysis, revealing several critical patterns. 
First, parallel implementations consistently outperform their sequential counterparts, with the acceleration factor scaling dramatically with graph size. 
For the Physics dataset, the largest in our evaluation, parallel training alone provides a 3.23$\times$ speedup over sequential LineML. 
When combined with pruning, this acceleration compounds substantially.

Most notably, \textbf{Parallel kNN Pruning emerges as the optimal configuration}, achieving a remarkable 13.98$\times$ speedup on Physics.
Also, kNN Pruning is capable of maintaining strong performance on smaller datasets (4.30$\times$ on Cora, 7.38$\times$ on CiteSeer, 3.99$\times$ on PubMed). 
This represents an effect where pruning reduces computational load while parallelization distributes the remaining work efficiently. 
The degree of acceleration correlates strongly with graph size: on Physics, parallel kNN pruning is 2.7$\times$ faster than the next best method (parallel degree pruning at 5.22$\times$), while on PubMed the advantage is more modest (1.26$\times$ faster than parallel degree pruning).

\begin{figure}[!ht]
    \centering
    \includegraphics[width=0.95\textwidth]{figs/Fused_All_Pruning_Comparison.pdf}
    \caption{Comparative analysis of training time speedup across pruning methods and datasets.}
    \label{fig:fused_pruning_comparison}
\end{figure}


The inference time results reveal a more nuanced trade-off between parallelization overhead and computational efficiency. 
For small to medium-sized graphs (Cora and CiteSeer), parallel implementations exhibit increased inference latency compared to their sequential counterparts. 
This overhead stems from inter-process communication costs that outweigh the benefits of parallel computation when graph partitions are small. 
Specifically, on Cora, parallel methods show 3.6-5.3$\times$ slowdown in inference time relative to sequential baselines.

However, this pattern reverses for larger graphs. 
On PubMed and Physics datasets, parallel kNN pruning achieves inference speedups of 1.71$\times$ and 4.24$\times$ respectively, demonstrating that the communication overhead becomes negligible relative to the computational savings when processing millions of edges. 
This crossover point, where parallelization becomes beneficial for inference, occurs at approximately 50,000 edges in the line graph, providing practical guidance for deployment decisions.


The most significant finding concerns the scaling behavior of different pruning strategies under parallel execution. 
As shown in Figure~\ref{fig:fused_pruning_comparison}, kNN pruning demonstrates near-ideal scaling: its parallel implementation achieves speedups that are multiplicative rather than additive relative to the sequential pruned version. 
For instance, on Physics, sequential kNN pruning provides a 4.60$\times$ speedup over baseline LineML, while parallel kNN pruning achieves 13.98$\times$—approximately the product of individual improvements (4.60$\times$ from pruning $\times$ 3.23$\times$ from parallelization $\approx$ 14.86$\times$). 
This indicates minimal parallelization overhead when combined with kNN pruning's regular, feature-based sparsity pattern.

Based on our HPC analysis, kNN pruning shows strong synergy with parallel architectures by improving communication efficiency and maintaining balanced workloads through its uniform neighborhood structure. We observe a clear scale-dependent behavior. For small graphs (under 50,000 edges), standard graph processing is sufficient: pruning provides only marginal gains, and parallel execution introduces unnecessary overhead. For medium-scale graphs (50,000–500,000 edges), efficient pruning yields substantial benefits, while sequential training remains preferable, as parallelization still incurs noticeable overhead. In contrast, large-scale graphs (over 500,000 edges) achieve near-ideal speedup under fully distributed execution.




\subsubsection{Comparative Analysis Against State-of-the-Art Methods}
\label{subsec:baseline_comparison}

To contextualize the performance of our proposed framework, we conduct a comprehensive comparison against six state-of-the-art graph learning methods described in the section \ref{subsub:Basline}: LGLP, SEAL, GAE, GraphSAGE, EG-DNMF, and DGI.
For this comparison, we employ the optimal configuration of LineML identified in previous sections: specifically, parallel LineML with kNN pruning (k=10).
the results, presented in Figure~\ref{fig:fused_baseline_comparison}, reveals critical insights about the efficiency and scalability trade-offs between line graph approaches and methods operating on the original graph topology.

Our analysis demonstrates that Parallel LineML achieves substantial speedups over the baseline line graph method LGLP, particularly on large-scale datasets. 
On the Physics dataset, Parallel LineML trains 14.09$\times$ faster than LGLP.
This dramatic improvement stems from two key innovations: our kNN pruning reduces the line graph edge count by 75.9\%, while parallelization efficiently distributes the computational load.

When compared to methods operating on the original graph Parallel LineML shows competitive scaling. 
For example, on Physics, the original graph contains only 34,493 nodes and 247,962 edges, while its line graph representation comprises 991,848 nodes and approximately 45 million edges before pruning. 
Despite processing 181$\times$ more edges than methods like GraphSAGE, Parallel LineML achieves training times within an order of magnitude as see in Figures \ref{fig:fused_baseline_comparison}. 
This demonstrates that our pruning and parallelization strategies effectively bridge the computational gap between line graph and original graph methods.


Inference results show scale-dependent behavior. On small datasets (Cora and CiteSeer), Parallel LineML incurs higher latency than original-graph methods due to parallel communication overhead; for example, SEAL achieves 7.48$\times$ faster inference on Cora. 
This trend reverses on larger graphs. 
On Physics, Parallel LineML outperforms LGLP by 4.16$\times$ (0.708\,s vs.\ 2.948\,s), with per-node latency of 0.0143\,ms, representing a 4.16$\times$ improvement over LGLP, though still higher than DGI (0.00055\,ms) operating on the original graph.
Scaling analysis (Figure~\ref{fig:fused_baseline_comparison}) further shows that the performance gap between line graph and original-graph methods widens with graph size: on Cora, GAE trains 12.51$\times$ faster than LGLP, while on Physics, DGI is 172.23$\times$ faster. Crucially, Parallel LineML alleviates this limitation through kNN pruning. While LGLP becomes prohibitively expensive, Parallel LineML preserves practical training times by reducing the effective graph size. On Physics, pruning removes 74.9\% of edges (from 45M to 10.8M), yielding a 14.09$\times$ speedup while maintaining link prediction accuracy.


This comparative analysis reveals fundamental trade-offs between expressive power and computational efficiency. 
Line graph methods, offer enhanced representational capacity for link prediction by explicitly modeling edge relationships but incur higher computational costs due to larger graph sizes. 
Original graph methods achieve superior efficiency by operating on the smaller primal graph but may sacrifice expressiveness for certain relational tasks.


\begin{figure}[t!]
    \centering
    \includegraphics[width=\linewidth]{figs/Fused_All_Baseline_Comparison.pdf}
    \caption{Performance comparison of LineML against baseline methods across four datasets. }
    \label{fig:fused_baseline_comparison}
\end{figure}

\section{Systematic Ablation Analysis and Hyperparameter Optimization}
\label{sec:SysAblation}
This section presents a comprehensive ablation study to identify the key architectural and optimization components driving LineML's performance.
The analysis is organized into three complementary dimensions: 
(1) architectural foundations of the GNN backbone, 
(2) the metric learning framework and sampling strategies, and 
(3) hyperparameter sensitivity and model capacity. 


For detailed ablation studies on hyperparameter sensitivity, metric learning components, and core architectural choices, we focus on a representative subset of seven datasets: BUP, UAL, GRQ, SMG, PubMed, YST, and CEG. 
This subset is carefully chosen to reflect variations in domain, size, and density, enabling a focused yet comprehensive analysis of how different design decisions influence performance across diverse graph characteristics.

\subsection{Architectural Foundations: GNN Backbones and Aggregation Strategies}

This subsection evaluates the core architectural choices for processing line graphs. 
We first compare GNN backbones (GraphSAGE vs. GCN vs. GAT) to identify the optimal foundation. 
We then analyze aggregation functions (mean vs. min vs. max) within the best-performing architecture.

\subsubsection{GNN Architecture Comparison}
\label{subsubsec:gnn_arch_comparison}
Figure~\ref{fig:gnn_arch_aggregator}a compares three prominent architectures. 
GraphSAGE achieves superior performance with a median ROC-AUC of 0.970 and minimal variability (IQR=0.012), while GAT (0.859 AUC, IQR=0.054) and GCN (0.873 AUC, IQR=0.064) show both lower accuracy and greater dataset dependence. 
This finding contrasts with some prior work on primal graphs where GAT often excels, highlighting that architectural preferences differ substantially between primal and line graph domains.

The performance gap stems from line graphs' structural characteristics. 
Line graphs transform each edge into a node, creating a new graph where node degrees scale with the number of triangles in the original graph. 
This produces dense neighborhoods where full neighborhood aggregation GCN or pairwise attention computation (GAT) becomes both computationally expensive and prone to over-smoothing. 
GraphSAGE's fixed-size neighborhood sampling mitigates both issues: it bounds computational cost while introducing stochasticity that regularizes against feature homogenization. 
This advantage is particularly pronounced in large-scale graphs where line graph nodes can reach thousands of neighbors.


\subsubsection{Aggregation Function Analysis}
\label{subsubsec:aggregation_analysis}

We further investigate how neighborhood information should be aggregated in GraphSAGE. 
The choice between min, max, and mean operators presents a trade-off between preserving distinctive features (min/max) and computing representative averages (mean). 
Figure~\ref{fig:gnn_arch_aggregator}b reveals that the min aggregator achieves the highest median AUC (0.967) with strong consistency (IQR=0.020), while mean performs worst (0.900 AUC, IQR=0.032).

This result aligns with principles from graph signal processing: the min and max operators act as non-linear filters that preserve extreme feature values, which often correspond to discriminative structural patterns in link prediction. 
In contrast, averaging (mean) performs linear smoothing that can dilute distinctive signals a particularly detrimental effect in line graphs where neighborhood information redundancy is high. 
The max aggregator (0.945 AUC, IQR=0.017) shows similar benefits, suggesting that selective, non-linear aggregation generally outperforms averaging for this task.

\begin{figure}[!ht]
\centering
\begin{subfigure}[b]{0.48\textwidth}
    \centering
    \includegraphics[width=\linewidth]{figs/1-gnn_type_ROC_boxplot.pdf}
    \caption{Comparison of GNN architectures.}
    \label{fig:gnn_architectures}
\end{subfigure}
\hfill
\begin{subfigure}[b]{0.48\textwidth}
    \centering
    \includegraphics[width=\linewidth]{figs/1-gnn_aggr_ROC_boxplot.pdf}
    \caption{Comparison of GraphSAGE aggregators.}
    \label{fig:graphsage_aggregators}
\end{subfigure}
\caption{Performance comparison of different GNN architectures and GraphSAGE aggregators for link prediction.}
\label{fig:gnn_arch_aggregator}
\end{figure}



\subsection{Metric Learning Framework Analysis}
\label{subsec:metric_learning_analysis}

We next investigate how metric learning components interact to shape the embedding space. 
We examine the sensitivity of key hyperparameters and compare alternative metric learning objectives.
Then, the central question is: \textit{How should negative sampling strategies and loss functions be coordinated to learn effective link representations?}


\subsubsection{Effect of Metric Learning Hyperparameters}
\label{subsubsec:metric_hyperparams}

The adaptive triplet loss in LineML is controlled by two main hyperparameters: the base margin $\alpha$ and the adaptive scaling factor $\beta$. 
Figure~\ref{fig:metric_hyperparameters} summarizes their influence on performance across multiple datasets.

\begin{figure}[!ht]
\centering
\begin{subfigure}[b]{0.48\textwidth}
    \centering
    \includegraphics[width=\linewidth]{figs/1-alpha_ROC_boxplot.pdf}
    \caption{Base margin $\alpha$}
    \label{fig:alpha_analysis}
\end{subfigure}
\hfill
\begin{subfigure}[b]{0.48\textwidth}
    \centering
    \includegraphics[width=\linewidth]{figs/beta_boxplot.pdf}
    \caption{Adaptive scaling $\beta$}
    \label{fig:beta_analysis}
\end{subfigure}
\caption{Sensitivity of metric learning hyperparameters.}
\label{fig:metric_hyperparameters}
\end{figure}

As shown in Figure~\ref{fig:alpha_analysis}, performance follows an inverted U-shaped trend, peaking at $\alpha=1.0$. 
Smaller values ($\alpha \leq 0.5$) provide insufficient separation, whereas larger values ($\alpha \geq 1.5$) impose overly strict constraints that hinder optimization. 
The narrow interquartile range at $\alpha=1.0$ indicates stable behavior across datasets, making it a reliable default choice.

The adaptive scaling factor $\beta$ controls the emphasis placed on difficult triplets during training. 
Figure~\ref{fig:beta_analysis} shows a clear optimum around $\beta=1.0$.
This setting yields a twofold gradient amplification (Proposition~\ref{prop:gradient_properties}) while increasing the number of active triplets by approximately 50\%. 
Larger values ($\beta \geq 1.5$) introduce instability, reflected in higher variance, whereas smaller values ($\beta \leq 0.5$) provide insufficient learning signal for challenging examples.

Overall, $\alpha$ governs the absolute separation in the embedding space, while $\beta$ determines the relative focus on hard examples. 
The combination $(\alpha=1.0, \beta=1.0)$ offers an effective balance, enabling meaningful separation without overemphasizing outliers. 
This configuration remains robust across all dataset categories, from dense social graphs to sparse biological networks.

\subsubsection{Comparison of Metric Learning Approaches}
\label{subsubsec:metric_approaches}

We next compare different metric learning objectives to assess their impact on link prediction. 
Specifically, we evaluate adaptive triplet loss ``AD\_triplet'', standard triplet loss ``triplet'', cosine embedding loss ``cosine'', multi-similarity loss ``multi-similarity'', and a baseline trained without metric learning using cross-entropy alone. 
Figure~\ref{fig:metric_comparison} reports the corresponding ROC-AUC distributions.

The adaptive triplet loss achieves the highest median ROC-AUC with the smallest interquartile range, indicating both strong performance and high consistency. 
This advantage arises from its direct optimization of relative distances combined with adaptive scaling that strengthens learning on difficult triplets.
Standard triplet loss attains comparable accuracy but shows slightly higher variability, suggesting reduced robustness to differences in graph structure.
Cosine embedding loss provides moderate results with substantially higher variability, reflecting sensitivity to absolute embedding magnitudes rather than relative geometry.
Multi-similarity loss yields the lowest performance among metric learning methods. Although designed to capture complex similarity patterns, its added complexity appears to reduce stability in link prediction settings.

The superior performance of adaptive triplet loss can be attributed to three factors: (1) {direct distance optimization}, which aligns embedding geometry with the link prediction objective; (2) {adaptive learning dynamics}, where gradients are scaled according to triplet difficulty (Equation~\ref{eq:adaptive_triplet}); and (3) {computational efficiency}, enabled by online semi-hard negative sampling. Together, these properties make adaptive triplet loss both accurate and practical for large-scale deployment.

\begin{figure}[!ht]
  \centering
  \includegraphics[width=0.60\linewidth]{figs/1-Metric_loss_ROC_boxplot.pdf}
  \caption{ROC-AUC distributions for different metric learning approaches: adaptive triplet, standard triplet, cosine embedding, multi-similarity, and cross-entropy.}
  \label{fig:metric_comparison}
\end{figure}

\subsubsection{Impact of Negative Link Sampling Strategies}
\label{subsubsec:sampling_strategies}


Negative sampling is not merely a computational convenience but a critical pedagogical component that determines what the model learns. 
We investigate four strategies (Section~\ref{subsec:link_sampling}) to answer: \textit{Which negative sampling approach produces the most informative training signal for link prediction?}
Figure~\ref{fig:sampling_strategy} reports ROC-AUC distributions across all datasets, both with and without metric learning.

\begin{figure}[t!]
    \centering
    \includegraphics[width=0.8\linewidth]{figs/sampling_strategy_ROC_boxplot.pdf}
    \caption{ROC-AUC comparison of negative sampling strategies with and without metric learning.}
    \label{fig:sampling_strategy}
\end{figure}

Interestingly, random negative sampling combined with metric learning attains the highest median ROC-AUC. 
In contrast, degree-based negative sampling offers the strongest overall trade-off between performance and stability. 
With metric learning, it achieves a median ROC-AUC of $0.959$, and remains consistently competitive even without metric learning. 
As theoretically motivated in Section~\ref{subsec:link_sampling}, sampling negatives proportionally to degree products preserves network heterogeneity and emphasizes statistically plausible but absent edges. 
Common neighbor sampling yields lower performance and limited robustness. 
While effective in social networks dominated by triadic closure, this strategy introduces strong locality bias and does not generalize well to biological or infrastructure graphs.
Hard negative sampling performs worst in our setting and incurs substantial computational overhead. 
Although designed to generate challenging examples, its adaptive nature amplifies noise during early training and requires repeated similarity computations.
As an outcome, this leads to unstable optimization and poor scalability on large heterogeneous networks.
Metric learning consistently improves performance and reduces variability across all sampling strategies, most significantly for random sampling by imposing structured distance constraints. However, it cannot fully compensate for the lack of informative negatives inherent in naive approaches.
We therefore select degree-based negative sampling as LineML's default.
When paired with metric learning, it provides robust accuracy, low variance, and computational efficiency ideal for scalable deployment in high-performance computing settings.


\subsubsection{Impact of Negative Sampling Ratio}
\label{subsubsec:sampling_ratio}

We finally study the influence of the negative-to-positive sampling ratio, varying it from 1:1 to 100:1. Figure~\ref{fig:sampling_ratio} reports AUC, Average Precision (AP), and F1 scores.

\begin{figure}[t!]
    \centering
    \includegraphics[width=\linewidth]{figs/metrics_boxplots_grid.pdf}
    \caption{Effect of negative sampling ratio on AUC, AP, and F1. }
    \label{fig:sampling_ratio}
\end{figure}

AUC remains largely stable across all ratios, reflecting its intrinsic robustness to class imbalance. 
For this reason, we additionally report AP and F1, which are more sensitive to sparsity and threshold-dependent decision quality. 
Both metrics peak at a 2:1 negative-to-positive ratio and subsequently degrade as the ratio increases, with AP and F1 dropping by 27.6\% and 31.9\%, respectively, when moving from 2:1 to 100:1. 
This trend indicates that excessive negative sampling amplifies sparsity and overwhelms the positive signal, making accurate discrimination increasingly difficult.

Metric learning consistently improves performance across all ratios and becomes increasingly critical as imbalance grows. 
At higher ratios, it substantially mitigates the degradation of precision-based metrics, yielding gains of 42.1\% in AP and 48.84\% in F1. 
These improvements stem from the explicit enforcement of structured distances in the embedding space, which counteracts the sparsity induced by abundant negatives and preserves meaningful separation between positive and negative links. 
Overall, these results demonstrate that metric learning plays a central role in sustaining robust performance under increasing imbalance, effectively compensating for the loss of positive density as the sampling ratio rises.


\subsection{Hyperparameter Sensitivity and Model Capacity Analysis}
\label{subsec:hyperparameter_capacity}
The final analysis addresses practical deployment considerations: \textit{How should hyperparameters and model capacity be configured for optimal performance-efficiency trade-offs?} We examine learning dynamics, batch size effects, and architectural depth to establish guidelines for real-world implementation.

\subsubsection{Hyperparameter Sensitivity Analysis}
\label{subsubsec:hyperparameter_sensitivity}

To systematically evaluate the sensitivity of LineML to key hyperparameters, we conducted a comprehensive one-at-a-time sensitivity analysis. 
This methodology isolates each parameter's effect by varying it across a defined range while maintaining others at their optimal baseline values. 
As shown in Figure~\ref{fig:hyperparameter_sensitivity}, three critical training parameters exhibit distinct optimal regions that balance convergence speed, training stability, and final performance.

The GNN encoder learning rate demonstrates an inverted U-shaped relationship with performance (Figure~\ref{fig:gnn_lr_boxplot}), with optimal results at $0.001$. 
Lower rates ($0.0001$) lead to slow convergence and potential underfitting, while higher rates ($0.01$) cause training instability and degraded performance.
This pattern reflects the complex optimization landscape of GraphSAGE operating on line graphs, where appropriate gradient magnitudes are essential for navigating the high-dimensional parameter space.

In contrast, the link classifier learning rate requires much more conservative values (Figure~\ref{fig:lc_lr_boxplot}), with $10^{-6}$ achieving optimal performance. 
This two-order-of-magnitude difference from the GNN learning rate stems from the classifier's simpler architecture and its position at the end of the processing pipeline. 
Larger learning rates ($10^{-4}$) cause the classifier to overfit to training patterns.
While the smallest tested value ($10^{-6}$) enables precise weight adjustments that preserve the discriminative information encoded by the GNN and metric learning components.

Batch size analysis reveals a clear trade-off between gradient estimation quality and optimization efficiency (Figure~\ref{fig:batch_size_boxplot}). Performance increases with batch size up to an optimal point of $256$, beyond which additional gains are minimal or negative. 
This pattern aligns with theoretical expectations: larger batches provide more stable gradient estimates but may converge to sharp minima with poorer generalization. 
The optimal batch size of $256$ represents a balance suitable for HPC environments, where memory constraints and parallelization efficiency must be considered alongside accuracy.

\begin{figure}[!ht]
\centering
\begin{subfigure}[b]{0.32\textwidth}
    \centering
    \includegraphics[width=\linewidth]{figs/1-gnn_lr_ROC_boxplot.pdf}
    \caption{GNN learning rate}
    \label{fig:gnn_lr_boxplot}
\end{subfigure}
\hfill
\begin{subfigure}[b]{0.32\textwidth}
    \centering
    \includegraphics[width=\linewidth]{figs/1-lr_lc_ROC_boxplot.pdf}
    \caption{Classifier learning rate}
    \label{fig:lc_lr_boxplot}
\end{subfigure}
\hfill
\begin{subfigure}[b]{0.32\textwidth}
    \centering
    \includegraphics[width=\linewidth]{figs/1-batch_size_ROC_boxplot.pdf}
    \caption{Batch size}
    \label{fig:batch_size_boxplot}
\end{subfigure}
\caption{Training hyperparameter sensitivity analysis. }
\label{fig:hyperparameter_sensitivity}
\end{figure}

\subsubsection{Model Capacity Analysis: Depth and Ensemble Size}
\label{subsubsec:model_capacity}

Model capacity is a central design factor that balances expressivity, generalization, and computational cost. 
We analyze two complementary dimensions: architectural depth and ensemble width. 
Figure~\ref{fig:model_capacity} shows clear optimal regions, beyond which additional capacity yields diminishing returns or degrades performance.

The interaction between GNN depth and classifier depth (Figure~\ref{fig:gnn_classifier_depth}) follows a non-monotonic trend, with peak performance at $3$ GNN layers and $2$ classifier layers. 
Shallow configurations (1--2 layers) underfit and fail to capture higher-order structural dependencies. 
In contrast, deeper models (4--5 layers) suffer from over-smoothing, an effect amplified in line graphs where high connectivity accelerates feature mixing across multiple hops. 
The optimal configuration balances these effects: three GNN layers capture sufficient neighborhood context, while two classifier layers preserve discriminative capacity without introducing unnecessary complexity.

The impact of ensemble size is illustrated in Figure~\ref{fig:classifier_heads}. 
Increasing the number of heads from $1$ to $4$ produces substantial gains, reflecting effective variance reduction through model diversity. 
However, expanding to $8$ heads provides only marginal improvement while nearly doubling computational cost. 
This behavior indicates that additional heads become highly correlated, limiting further benefits.


Moderate depth combined with a small ensemble achieves strong performance while avoiding over-smoothing and excessive computation. 
Based on this analysis, we adopt $3$ GNN layers, $2$ classifier layers, and $4$ ensemble heads as the default configuration, offering an effective balance between accuracy, stability, and scalability.


\begin{figure}[!ht]
\centering
\begin{subfigure}[b]{0.48\textwidth}
    \centering
    \includegraphics[width=\linewidth]{figs/gnn_classifier_heatmap.pdf}
    \caption{GNN-classifier depth interaction}
    \label{fig:gnn_classifier_depth}
\end{subfigure}
\hfill
\begin{subfigure}[b]{0.48\textwidth}
    \centering
    \includegraphics[width=\linewidth]{figs/classifier_heads_boxplot.pdf}
    \caption{Classifier ensemble size}
    \label{fig:classifier_heads}
\end{subfigure}
\caption{Model capacity optimization. (a) Depth interaction heatmap reveals optimal configuration at 3 GNN layers and 2 classifier layers, with deeper models suffering from over-smoothing. (b) Ensemble size analysis shows diminishing returns beyond 4 heads.}
\label{fig:model_capacity}
\end{figure}

\section{Discussion and Limitations}
\label{sec:discussion}

\subsection{Discussion}


The experimental results show that LineML consistently achieves state-of-the-art performance across diverse network types, validating its theoretical design. 
This section interprets these outcomes through the framework’s three core components—line graph reformulation, adaptive metric learning, and principled pruning—and explains how their interaction enables robust link prediction.

LineML derives its strongest gains from transforming links into first-class entities via the line graph. 
As shown in Theorem~\ref{thm:wl_expressivity}, a GNN operating on $L(G)$ matches the expressive power of the 1-WL test for edge discrimination. 
This property translates empirically into superior performance on heterogeneous and sparse networks such as ADV (AUC: 95.51\%), BUP (98.78\%), HPD (96.48\%), and CEG (96.01\%). 
In these settings, links encode complex, non-homophilic relations that node-centric models fail to capture effectively. 
Performance differences across datasets further support this interpretation. 
On dense and homogeneous citation networks (e.g., NSC), the advantage of the line graph is marginal (99.82\% for LGLP vs.\ 99.51\% for LineML). 
In contrast, on large and structurally complex networks such as Physics, LineML becomes clearly superior (97.70\% vs.\ 97.25\%), indicating that the representational benefit of $L(G)$ increases with network heterogeneity and scale.

Adaptive metric learning plays a central role in shaping a task-aware embedding space. 
Proposition~\ref{prop:gradient_properties} shows that the adaptive triplet loss amplifies gradients on semi-hard triplets, improving robustness under class imbalance. 
Ablation results (Figure~\ref{fig:metric_comparison}) confirm this effect, with the adaptive loss achieving the most stable median AUC (0.970) and a narrow interquartile range of 0.012. 
Its impact is strongest when combined with degree-based negative sampling, which provides structurally plausible negatives. 
Hyperparameter analysis (Figure~\ref{fig:metric_hyperparameters}) identifies $(\alpha=1.0,\beta=1.0)$ as the optimal configuration, balancing margin enforcement and adaptive focus. 
This setting yields consistent performance across datasets, even with limited training data (30\%--40\% in Figure~\ref{fig:training_size_effect}). 
As class imbalance increases, the benefit of metric learning becomes more pronounced, improving Average Precision by up to 42.1\%.

Pruning contributes both scalability and predictive accuracy. 
Theoretical analysis frames pruning as spectral approximation, and experiments confirm that kNN-based pruning ($k=10$) removes approximately 76\% of line graph edges while preserving or improving performance (e.g., +0.38\% AUC on Physics). 
This behavior aligns with sparsification theory: removing redundant or noisy connections acts as denoising while retaining geometrically meaningful neighborhoods. 
kNN pruning consistently outperforms degree and spectrum alternatives, achieving higher AUC with greater edge reduction. 
Parallel execution further amplifies these gains, delivering near-linear speedups (13.98$\times$ on Physics) and validating Lemma~\ref{lemma:trainingTimeReduction}. 
On large graphs, pruning becomes essential, reducing training time from 1,793\,s to 389\,s per epoch on Physics while maintaining 98.07\% AUC.




\subsection{Limitations}
\label{sec:limitations}


Despite its strong performance, LineML has several limitations that point to fruitful research directions. 
A primary limitation of LineML lies in the cost of explicit line graph construction. 
Although Theorem~\ref{thm:wl_expressivity} demonstrates its expressive advantages, building $L(G)$ requires $O(\sum_v d_v^2)$ time and memory. 
This becomes prohibitive for graphs with highly skewed degree distributions, where a small number of hubs can cause the line graph to grow super-linearly with $|V|$. In extreme cases, such as web-scale networks, $L(G)$ may contain billions of edges, exceeding practical memory limits even on HPC systems.
While our pruning strategies substantially reduce graph size, they operate after the full line graph has been constructed and therefore do not eliminate the initial memory overhead. As a result, LineML is not designed for ultra-large graphs, or deployment on memory-constrained devices. These settings are outside the scope of the current work.


The theoretical guarantees underpinning LineML assume  static graphs with discrete node labels. 
This represents a significant limitation when applied to real-world networks that are often directed, temporal, and richly attributed. 
For temporal graphs, Theorem~\ref{thm:leakage_free} addresses leakage prevention but does not provide guidance on how to optimally incorporate temporal edge features or how the line graph's structure evolves over time. 
The current feature transformation (Hadamard product with gating in Eq.~\ref{eq:attribute_transform_with_edge}) is effective for continuous attributes but may be suboptimal for categorical, multi-modal, or missing features.


LineML's performance, like all contrastive learning frameworks, is inherently tied to the quality of negative sampling. 
While degree-based sampling proved robust across our benchmark datasets, it remains a heuristic choice that may fail in specialized domains. 
For instance, in bipartite graphs or graphs with strong community structure, degree-based sampling may not generate sufficiently challenging negatives. Moreover, the sampling distribution is static throughout training, whereas an adaptive curriculum that evolves with the model's learning progress could be more effective. 
Theoretical analysis of the optimal negative sampling distribution for line graph embeddings remains an open question. 

LineML is designed for simple graphs with binary edge existence. 
Many real-world relational systems involve multiple edge types (e.g., social networks with friend, family, colleague relations) or higher-order interactions represented as hyperedges (e.g., co-authorship with multiple authors, chemical reactions with multiple reactants). 
The current framework does not naturally extend to these settings. 
Converting a multi-relational graph to a line graph would require deciding how to handle different relation types in the adjacency construction, potentially losing important semantic distinctions. 
For hypergraphs, the line graph definition becomes even more complex, as hyperedges can intersect in various patterns beyond simple node sharing. 










\section{Conclusion and Future Work}
\label{sec:conclusion}

This paper addressed the link prediction problem, a fundamental task in network analysis that remains challenging due to indirect edge modeling and severe class imbalance in real-world graphs. 
While reformulating link prediction as node classification on line graphs offers a principled solution by treating edges as first-class entities, this approach introduces major computational challenges caused by the quadratic growth of the transformed graph. 
To overcome these limitations, we proposed \textbf{LineML}, an end-to-end framework that converts link prediction into a node classification formulation on line graphs while remaining scalable in practice. 
LineML integrates a GraphSAGE-based encoder to capture structural and attribute information, an adaptive metric learning module with degree-biased negative sampling to handle class imbalance, and a pruning and parallel training system to mitigate quadratic complexity. 
Together, these components enable direct edge modeling and robust representation learning under challenging data distributions.
Extensive experiments conducted on 18 benchmark datasets demonstrated the effectiveness of the proposed framework. 
LineML achieved the highest average rank on social and biological networks (1.0) and the best average rank on citation networks (2.4). 
Statistical analysis confirmed near-complete dominance on social and biological datasets (Cliff’s $\delta \geq 0.97$ against all 14 baselines) and substantial improvements on citation networks (Cliff’s $\delta \geq 0.73$ against 12 of 14 methods). 
All observed improvements were statistically significant according to the Wilcoxon signed-rank test ($p < 0.01$).
From a computational perspective, this work demonstrated that high-performance computing is essential for scalable line graph learning. 
Spectral pruning reduced training time by up to 4.6$\times$ while preserving predictive performance.
The distributed pruning implementation achieved near-linear scaling, delivering a 13.98$\times$ speedup over line graph baselines.
These results confirm that pruning and parallel training are key enablers for applying line graph-based models to large networks.

Future work will pursue several promising directions. 
First, we plan to extend the theoretical framework to temporal, and multi-relational graphs, enabling LineML to model dynamic link evolution. 
Second, developing uncertainty quantification for line graph embeddings is essential for deployment in risk-sensitive applications, where calibrated confidence estimates are required alongside predictions. 
Third, we will explore adaptive negative sampling strategies that evolve during training, allowing the model to progressively focus on more informative hard negatives while maintaining optimization stability. 
Finally, integrating automated hyperparameter optimization will improve usability and reproducibility, reducing manual tuning and facilitating broader adoption.










\section*{Funding}
This research received no external funding.

\section*{Conflict of Interest}
The authors declare that they have no conflicts of interest.

\section*{Data Availability}
The datasets used in this study are publicly available online.

\section*{Code Availability}
The code used for this research is available upon request.

\section*{Author Contributions}
Adil Imad Eddine Hosni conceptualized the study, designed the methodology, and wrote the initial manuscript draft.  
Islam Baira contributed to data preprocessing, analysis, and manuscript revisions.  
Badis Djamaa was responsible for implementing the experiments and optimizing model performance.  
M’hamed Mataoui supervised the research, reviewed the manuscript, and provided critical feedback.  
Sidhoum Abdellah Hamouda provided theoretical insights and assisted in manuscript structuring.  
Hichem Merini worked on data preprocessing, analysis, and manuscript revisions.  
Mohamed Chakib Amrani contributed to data collection, literature review, and statistical analysis.  
All authors read and approved the final manuscript.

\section*{Declaration of Generative AI and AI-Assisted Technologies in the Writing Process}

During the preparation of this work, the authors used AI-assisted tools solely for language proofreading. The ideas, experimental protocol, and manuscript writing were entirely conceived, developed, and executed by the authors. After using AI for language refinement, the authors carefully reviewed and edited the content as needed and take full responsibility for the final version of the publication.

\bibliography{sn-bibliography}



% \section*{Declarations}

% Some journals require declarations to be submitted in a standardised format. Please check the Instructions for Authors of the journal to which you are submitting to see if you need to complete this section. If yes, your manuscript must contain the following sections under the heading `Declarations':

% \begin{itemize}
% \item Funding
% \item Conflict of interest/Competing interests (check journal-specific guidelines for which heading to use)
% \item Ethics approval and consent to participate
% \item Consent for publication
% \item Data availability 
% \item Materials availability
% \item Code availability 
% \item Author contribution
% \end{itemize}



\end{document}
